\documentclass[12pt,a4paper]{article}

%% --- Encoding and fonts ---
\usepackage[utf8]{inputenc}
\usepackage[T1]{fontenc}
\usepackage{lmodern}

%% --- Language ---
\usepackage[english]{babel}

%% --- Mathematics ---
\usepackage{amsmath, amssymb, amsthm}
\usepackage{mathtools}

%% --- Layout ---
\usepackage[margin=1in]{geometry}
\usepackage{setspace}
\onehalfspacing

%% --- Figures and tables ---
\usepackage{graphicx}
\usepackage{booktabs}
\usepackage{caption}
\usepackage{subcaption}
\usepackage{float}

%% --- Appendix ---
\usepackage[title,titletoc]{appendix}

%% --- Cross-references and links ---
\usepackage{hyperref}
\hypersetup{colorlinks=true, linkcolor=blue, citecolor=blue, urlcolor=blue}

%% --- Bibliography ---
\usepackage{natbib}
\bibliographystyle{apalike}

%% --- Draft tools ---
\usepackage{xcolor}
\newcommand{\todo}[1]{\textcolor{red}{\textbf{[TODO: #1]}}}
\newcommand{\note}[1]{\textcolor{blue}{\textit{[Note: #1]}}}
% --- Custom Math Commands ---
\newcommand{\E}{\mathbb{E}}
\newcommand{\R}{\mathbb{R}}
\newcommand{\cA}{\mathcal{A}}          % asset space
\newcommand{\cH}{\mathcal{H}}          % productivity space
\newcommand{\cM}{\mathcal{M}}          % health space
\newcommand{\cF}{\mathcal{F}}          % formality space
\newcommand{\cQ}{\mathcal{Q}}          % pension-counter space
\newcommand{\ind}{\mathbf{1}}           % indicator function
\newcommand{\phij}{\phi_j}             % cohort weight
\newcommand{\jR}{j_R}                  % retirement age
% Type-specific shortcuts
\newcommand{\tL}{\theta_L}
\newcommand{\tH}{\theta_H}

%% --- Theorem environments ---
\newtheorem{assumption}{Assumption}
\newtheorem{proposition}{Proposition}
\newtheorem{definition}{Definition}
\newtheorem{remark}{Remark}

%% ====================================================================
\begin{document}

%% --- Title block ---
\title{\textbf{Spending Smarter under Demographic Pressure:}\\[6pt]
Fiscal Efficiency, Health, and Gender Dynamics\\
in Latin America\thanks{The work was financed with the support of the Latin America and the Caribbean Research Network of the Inter-American Development Bank. The opinions expressed in this publication are those of the authors and do not necessarily reflect the views of the Inter-American Development Bank, its Board of Directors, or the countries they represent.}}

\author{
Diego Ascarza-Mendoza\textsuperscript{1} \quad Hermilo Cort\'es\textsuperscript{1}\\[3pt]
Judith M\'endez\textsuperscript{1} \quad H\'ector J.\ Villarreal\textsuperscript{1}
\\[6pt]
\textsuperscript{1}Iniciativa para la Transici\'on Econ\'omica y Demogr\'afica (ITED),\\
Escuela de Gobierno y Transformaci\'on P\'ublica,\\
Tecnol\'ogico de Monterrey
}

\date{%
Prepared for the Inter-American Development Bank,\\
Fiscal Division\\[12pt]
\textit{Preliminary draft -- June 2026}\\
\textit{Please do not cite or circulate without permission}
}

\maketitle

%% --- Abstract ---
\begin{abstract}
Latin America is aging before it has grown rich, and the fiscal pressure this places on health and pension systems is sharpened by persistent labor informality and limited fiscal space. We develop a stochastic overlapping-generations model with endogenous health investment, agents heterogeneous in sex and skill, and a pay-as-you-go pension system whose contribution rate clears endogenously. Health capital depreciates with age, raises survival, productivity, and the amenity value of consumption, and is restored through medical spending, so it links longevity to the contribution rate through a fiscal amplification loop whose net sign is ambiguous. The model is designed for M\'exico, Costa Rica, and Panam\'a; this version reports the Mexican calibration, estimating the labor-efficiency profile and skill premium from ENOE micro-data and the remaining preference and health-technology parameters by simulated method of moments. Three findings follow. At M\'exico's projected 2050 demographics, the 2020 contribution rate and benefit formula open a pension financing gap that no finite stationary stock of public debt can close, narrowing the menu of feasible long-run responses below what the policy debate assumes. The calibrated capital stock is governed first by the persistence of idiosyncratic income risk and only second by skill composition. And a replacement-rate cut is regressive, while under unchanged policy the equilibrium contribution rate rises from 9.3\% to 15.9\% as the population ages.
\end{abstract}

\noindent\textbf{Keywords:} demographic transition; fiscal sustainability; health expenditure; overlapping generations; labor informality; Latin America

\noindent\textbf{JEL codes:} E62, H51, I13, J11, O17

%% ====================================================================
%% SECTION 1 — INTRODUCTION
%% ====================================================================
\section{Introduction}
\label{sec:intro}

Latin America is aging. Yet it is doing so under conditions that distinguish the region sharply from the experience of high-income countries: fertility decline is rapid, formal labor markets remain incomplete, fiscal space is tight, and public health systems are chronically underfunded. The conjunction of these forces creates a structural dilemma that is at once demographic, fiscal, and institutional. The central question this paper addresses is how countries in the region can reallocate public health resources---spending \textit{smarter}---to meet the growing burden of aging populations without compromising fiscal sustainability or weakening the social security institutions on which workers depend.

The demographic challenge is well documented but its fiscal dimensions remain incompletely understood. By 2045, old-age dependency ratios in M\'exico, Costa Rica, and Panam\'a---our three country cases---are projected to at least double relative to their 2020 levels \citep{ECLAC2022}. Unlike the transitions observed in Western Europe or East Asia, this aging is occurring before these societies have achieved high-income status---a phenomenon often described as \textit{growing old before growing rich}. The implication is that the fiscal base from which health and pension commitments must be financed is narrower, and the informal sector, which contributes little to social security revenues, remains large. Critically, the aging that took Europe roughly 70 years to accomplish will take Latin America approximately 25 years, leaving far less time to legislate and implement the fiscal adjustments that demographic change demands \citep{WorldBankInformalizing2025}.

M\'exico illustrates the challenge in its starkest form. Between 2015 and 2024, our estimates show that Mexico's fiscal health gap widened from 3.0 to 4.8 percent of GDP: actual public health expenditure, held at roughly 2.8 percent of GDP, now covers barely more than a third of demographically projected requirements. Costa Rica and Panam\'a present a more favorable picture, both countries currently spend approximately in line with estimated needs, but demographic projections imply rising pressures in both cases as population aging accelerates over the next two decades. The cross-country variation in fiscal gaps, institutional architectures, and informality profiles is precisely why a comparative analysis of three countries is essential: it reveals which features of health-system design are structurally protective and which are artifacts of a particular country's fiscal position at a given moment. M\'exico's combination of high informality and a deteriorating fiscal gap identifies the most severe case; Costa Rica's historically generous contributory system through the CCSS (\textit{Caja Costarricense de Seguro Social}) shows the trajectory of a system approaching the limits of its fiscal surplus in health; and Panam\'a's institutional fragmentation between the CSS and the Ministry of Health illustrates efficiency losses from dual-system design. Together, the three cases span the relevant policy space and allow us to distinguish between country-specific and structural determinants of fiscal sustainability under aging.

A distinctive feature of the Latin American demographic transition, one that separates it analytically from high-income country experience, is the \textit{premature health crossing point}. In M\'exico, the proportion of the population in poor health crosses 50 percent at approximately age 60, a decade earlier than the 70 to 75 age range observed in high-income countries. This implies that the health system faces peak expenditure pressure precisely at the moment when workers are exiting the labor force and the contribution base is contracting. Skill differentials compound this pattern: among low-skilled workers, the health crossing point occurs even earlier, with rapid deterioration beginning in the late 40s, while high-skilled workers maintain good health status on average into their late 60s. This 20-year gap in healthy longevity between skill groups is the central empirical fact around which our model is built. Health is not one dimension of fiscal pressure among many, it is the first-order driver of the sustainability problem, and it is the analytical center of this paper.

This paper's contribution is empirical and structural. We calibrate the model to M\'exico, substituting ENOE 2018 to 2024 labor-market micro-data for the sex-by-age efficiency profile, the skill premium, and the prime-age hours target, and estimating the remaining six preference and health-technology parameters by simulated method of moments; the calibration matches five of six target moments within two standard errors (three within one) and reveals a structural divergence on the value of a statistical life worth reporting on its own terms. Three quantitative findings follow. The strongest is structural: at M\'exico's projected 2050 demographic structure, the 2020 contribution rate combined with the 2020 benefit formula generates a pension financing gap that \textit{cannot be financed by a finite stationary stock of government debt}: the menu of feasible long-run policy responses to demographic aging in Latin America is structurally narrower than commonly assumed. Second, the calibrated baseline equilibrium holds substantially less capital than an illustrative literature-anchored parameterization; a Shapley decomposition attributes the larger share of the gap to the lower persistence of the empirically estimated income process (buffer-stock saving is weaker when shocks mean-revert) rather than to the low-skill-heavy cohort composition, so persistence governs the aggregate. Third, the policy experiments at the calibrated baseline confirm that a replacement-rate cut is regressive (low-skill types bear the loss while the high-skill gain) and that the equilibrium contribution rate rises by roughly 70 percent under the 2050 demographic structure with unchanged policy. Section~\ref{sec:calibration} reports the Mexican-data calibration and the calibrated baseline equilibrium, and Section~\ref{sec:experiments} the policy experiments.

We construct a stochastic overlapping-generations economy with endogenous health investment, heterogeneous agents who differ by skill type, and a government sector that operates a pay-as-you-go pension system. Health capital is a continuous state variable that depreciates with age and can be partially restored through medical spending, subject to diminishing returns. Health simultaneously affects survival probabilities, labor productivity, and the amenity value of consumption---making it a state variable that propagates through every margin of household decision-making. The model is solved as a stationary competitive equilibrium in which factor prices are determined endogenously through firm optimization and market clearing. The model is designed for M\'exico, Costa Rica, and Panam\'a, drawing on data from the Global Burden of Disease \citep{GBD2021}, the OECD Taxing Wages series \citep{OECDTaxingWages2026}, the WHO Global Health Expenditure Database \citep{WHOGHED2024}, and national household surveys; this version reports the Mexican calibration.

The model's health formation technology generates a \textit{longevity channel} that operates as a macroeconomic transmission mechanism \citep{Carrasco2026}. When households invest in health, they live longer, which strengthens the life-cycle savings motive, raises the aggregate capital stock, and alters the equilibrium interest rate and wage. This channel operates asymmetrically across skill groups: low-skilled workers, who face earlier and steeper health depreciation, have shorter effective planning horizons and accumulate less precautionary savings, reinforcing the health-income divergence that the premature health crossing point generates. In high-income countries, where health deterioration begins later and proceeds more gradually, the longevity channel operates more uniformly across the skill distribution. In Latin America, it amplifies inequality.

The longevity channel interacts with the pension system through a fiscal amplification mechanism. The pension contribution rate $\tau^p$ adjusts endogenously to balance the system each period, equaling the replacement rate $\kappa$ scaled by the dependency ratio $N^R/N^W$---the ratio of retirees to workers. As the population ages and the dependency ratio rises, $\tau^p$ increases mechanically, widening the tax wedge on formal labor income. Health investment that extends productive working lives partially offsets this pressure by keeping workers in the labor force longer and by improving their productivity while active. But the offset is partial: healthier workers who survive into retirement also draw pension benefits for more periods, raising $N^R$. The net fiscal effect of improved population health is therefore ambiguous in sign and quantitatively important, precisely the kind of question that requires a calibrated general-equilibrium model to resolve.

The model's structure generates several classes of quantitative results. First, lifecycle profiles of health capital, consumption, and savings by skill type, where the interaction of age-dependent health depreciation with endogenous investment produces divergent trajectories: high-skilled households maintain health longer, save more through the longevity channel, and arrive at retirement with substantially larger asset buffers. Second, the fiscal gap decomposition, applying age-expenditure profiles to the model's stationary population distribution yields country-specific estimates of the divergence between demographically projected health needs and actual public spending. Third, the equilibrium pension contribution rate responds endogenously to the dependency ratio, feeding back into household savings and labor supply decisions. Fourth, unlike partial-equilibrium analysis, this model allows health-driven changes in aggregate savings to move the interest rate and the wage, so that a deterioration in population health affects all agents through factor prices, not only the households directly hit. Fifth, the government budget closure ties the fiscal accounts together in general equilibrium, so that the sustainability of pension and health commitments can be read off the equilibrium contribution rate and the structure of the budget rather than assumed: a direct fiscal-sustainability diagnostic that identifies where adjustment pressures will be most severe.

The baseline model is calibrated to the formal-sector population. To study the fiscal consequences of expanding the contributory base, we design a formalization experiment. The total population $A$ is partitioned into formal workers $B$ and informal workers $C$, with high-skilled workers overrepresented in $B$ and low-skilled workers overrepresented in $C$. The experiment simulates a formalization program that incorporates a fraction of $C$ into $B$, using observed income and health profiles from national household surveys, ENOE for M\'exico, ENAHO for Costa Rica, and the Encuesta de Hogares for Panam\'a. The skill composition of the newly formalized cohort, predominantly low-skilled, with earlier health depreciation and lower lifetime earnings, generates competing fiscal effects: a broader contribution base that reduces the equilibrium pension rate $\tau^p$, against higher per-capita health expenditure claims and lower average productivity. The model quantifies the net fiscal impact and identifies, for each country, the threshold formalization rate at which the contribution-base effect dominates the expenditure-cost effect.

Beyond the formalization experiment, we examine a battery of policy counterfactuals that vary the fiscal instruments and institutional design parameters of the health and pension systems. These include: (i)~expansion of public health coverage through contributory versus tax-financed channels, with explicit tracking of effects on the tax wedge and the equilibrium contribution rate; (ii)~preventive health interventions targeted at ages 25--35, where the model's lifecycle profiles suggest the welfare return is steepest; (iii)~modified pension vesting requirements that alter the cumulative contribution threshold; and (iv)~joint experiments that combine health and pension reforms to identify complementarities. The joint experiments are essential because these instruments interact: a reduction in the pension vesting threshold raises the value of formal employment for workers most likely to be deterred by high contribution rates, and this effect is amplified when combined with preventive health investments that keep those workers productive longer. Both independent and joint counterfactuals are reported, addressing directly the question of whether policy complementarities are quantitatively important.

The model's architecture supports a natural extension toward gender-differentiated analysis. Rather than constructing a separate household bargaining module, we exploit the fact that sex-differentiated income profiles and survival probabilities already embed the cumulative effects of labor market discrimination, career interruptions from child-bearing, gendered sorting into formal and informal employment, and differences in care responsibilities. By expanding the skill-type space from two types ($\theta_L$, $\theta_H$) to four ($\theta_L^M$, $\theta_H^M$, $\theta_L^F$, $\theta_H^F$), calibrated with sex-specific earnings trajectories and the survival differentials documented in our data section---where women outlive men by 6--7 years across all three countries---the model can decompose fiscal sustainability by sex within the existing equilibrium framework.\footnote{In the current draft the sex-by-age efficiency profile $e^g_j$---and the implied $\approx 26\%$ prime-age sex earnings differential---is estimated from ENOE 2018--2024 (Section~\ref{sec:calibration}), while the sex-specific survival schedule retains a provisional asymmetry pending calibration from UN WPP 2024 sex-specific life tables (the September deliverable). The welfare levels by sex reported in Section~\ref{sec:experiments} are reportable in direction; their precise magnitudes remain provisional in the sex-survival dimension.} This decomposition asks: what fraction of the pension system's fiscal pressure comes from women's longer survival (drawing benefits for more periods) versus their lower contribution base (reflecting lower lifetime earnings)? The health expenditure decomposition operates analogously---women's later health crossing point but longer post-retirement survival generates a fiscal profile distinct from men's earlier deterioration but shorter lifespan. This calibration-based approach to gender lets the data carry the information that a structural bargaining model would generate endogenously, at a fraction of the computational cost and with direct empirical grounding.

The paper is structured as follows. Section~\ref{sec:lit} reviews the relevant literature on health financing, OLG models of aging, the longevity channel, formalization incentives, and gender heterogeneity in lifecycle models. Section~\ref{sec:model} presents the theoretical framework---a stochastic OLG economy with endogenous health investment, heterogeneous agents, and a PAYG pension system---and establishes that it admits a well-behaved stationary equilibrium. Section~\ref{sec:calibration} is the empirical core: it reports the Mexican-data calibration---the estimated parameters and their fit, the value-of-statistical-life result, and the calibrated baseline equilibrium that the experiments perturb. Section~\ref{sec:experiments} develops the experimental program at the calibrated baseline---stationary policy counterfactuals, the cost of inaction under 2050 demographics, the non-existence result under debt-financed inaction, and the formalization experiment. Section~\ref{sec:conclusion} concludes by summarizing the model's analytical contributions and presenting the research plan for the remaining phases of the project.

%% ====================================================================
%% SECTION 2 — LITERATURE REVIEW
%% ====================================================================
\section{Literature Review}
\label{sec:lit}

This paper connects five strands of literature: the economics of health financing and social insurance, the quantitative OLG literature on health and aging, the emerging work on endogenous health investment and longevity as macroeconomic channels, the empirical literature on informality and formalization incentives in Latin America, and the nascent efforts to incorporate gender heterogeneity into lifecycle fiscal models.

%% --------------------------------------------------------------------
\subsection{Health Financing, Social Insurance, and the Tax Wedge}
\label{subsec:lit_health}

The economics of health insurance in developing economies builds on the foundational insight that health shocks are both large and poorly insurable through private markets, generating welfare losses that can be substantially mitigated by pooling arrangements \citep{Arrow1963}. In Latin America, the design question is not whether to pool health risk but through what institutional channel. Systems organized around contributory social security---the Bismarckian tradition---link coverage to formal employment; systems organized around general taxation---the Beveridgean tradition---link coverage to residency or citizenship. In practice, as emphasized by \citet{Levy2008}, most countries in the region operate hybrid systems in which both channels coexist, often with fragmentation, duplication, and unequal benefit packages.

\citet{Summers1989} showed that the welfare implications of mandated benefits depend critically on whether workers value those benefits enough to accept compensating wage reductions. When they do, mandated contributions are essentially priced into wages and generate no net distortion; when they do not---either because benefits are not credible or because workers discount future benefits heavily---contributions function as a tax on formal labor, raising the cost of compliance and reducing formalization. \citet{WorldBankInformalizing2025} formalize this logic through the concept of the Formalization Tax Rate (FTR), which captures the net fiscal cost of formal employment including both mandatory contributions and foregone non-contributory benefits. Their simulations for a range of LAC countries demonstrate that the FTR is highest for low-income workers, precisely the group with the weakest attachment to the formal sector. The FTR framework provides the conceptual bridge between our model's tax wedge---the gap between gross labor cost and net take-home pay for formal versus informal workers---and the observed persistence of informality across the earnings distribution.

The empirical measurement of the tax wedge is provided by the OECD's \textit{Taxing Wages} series. \citet{OECDTaxingWages2026} report that in 2025 the total tax wedge for a single worker earning the average wage in Mexico was 21.7 percent of labour costs---composed of 11.8 percent income tax, 1.4 percent employee social security contributions, and approximately 8.5 percent employer contributions---compared with an OECD average of 35.1 percent. Costa Rica's tax wedge, at 9.8 percent, is among the lowest in the OECD. Two features of the Mexican data are particularly relevant for our calibration. First, the 2021 Social Security Law reform introduced a graduated employer contribution schedule for discharge and old-age insurance (\textit{Cesant\'ia en Edad Avanzada y Vejez}) that scales from 3.15 percent at one minimum wage to 11.875 percent at 4.5 UMAs by 2030, creating a progressive structure in which the marginal cost of formalizing a worker depends on the wage level. Second, the OECD's new progressivity indicator---which measures how the marginal tax wedge varies from 50 to 500 percent of the average wage---places Mexico and Costa Rica at the bottom of the OECD ranking, confirming that the tax system's interaction with informality operates not only through the level of the wedge but through its incidence across the wage distribution.

A central insight from this literature, and one that directly motivates our model's structure, is the asymmetry between health and pension benefits in workers' valuations. Health coverage generates immediate, tangible protection against catastrophic expenditure shocks; pension benefits, by contrast, deliver returns far in the future and are subject to credibility discounts that grow with institutional opacity and long vesting periods. \citet{BaezaPackard2006} argue that health coverage is the social insurance benefit that workers learn most about and value most highly---making health financing design a first-order determinant of compliance behavior. This asymmetry implies that the institutional channel through which health coverage is provided---contributory or tax-financed---has direct consequences for formalization incentives that are larger, per unit of fiscal cost, than those of pension design.

%% --------------------------------------------------------------------
\subsection{OLG Models of Fiscal Sustainability, Health, and Aging}
\label{subsec:lit_olg}

The quantitative OLG literature provides the methodological foundation for this paper. Since \citet{AuerbachKotlikoff1987}, the workhorse framework for analyzing the intergenerational fiscal effects of demographic change has been the multi-period OLG model with a government sector. The health-augmented extensions most relevant to this paper include the work of \citet{DeNardi2010}, who show that uncertain medical expenses are a primary driver of late-life savings among older households in the United States, and \citet{KuhnFeichtinger2015}, who analyze optimal health and retirement choices over the life cycle in a framework where health capital depreciates with age and can be partially restored through investment.

The Grossman tradition of modeling health as a depreciable capital stock---in which individuals invest to maintain or restore a stock that deteriorates with age---provides the theoretical foundation for our health formation equation. \citet{DalgaardStrulik2014} extend this framework to show how optimal aging and mortality patterns emerge endogenously from the interaction of health investment decisions with biological depreciation, yielding a structural interpretation of the Preston curve. \citet{Fioroni2010} analyzes optimal savings and health spending jointly over the life cycle, demonstrating that the interaction between precautionary savings motives and health investment generates consumption-savings profiles that differ substantially from models without endogenous health. \citet{FrankovicWrzaczek2020} study how medical progress interacts with overlapping-generations dynamics, showing that improvements in health technology alter not only individual life-cycle profiles but also aggregate capital accumulation and factor prices through general-equilibrium channels. \citet{YewZhang2018} introduce longevity externalities into a lifecycle-dynastic model, demonstrating that health spending decisions by one generation affect the survival and savings behavior of subsequent generations---a mechanism that, while not directly present in our stationary-equilibrium framework, underscores the importance of modeling health as an endogenous state variable rather than an exogenous process.

For the Latin American context, \citet{CortesEtAl2024} develop a partial-equilibrium OLG model for M\'exico in which agents differ by skill and face stochastic health transitions, health-dependent productivity, and age-varying out-of-pocket expenditure profiles. Their results---consumption gaps of 40--70 percent between good and poor health groups, and a 20-year gap in healthy longevity by skill level---form the empirical backbone of the present paper. The primary extension in this work is the move to general equilibrium with endogenous wages, an explicit pension system with endogenous contribution rates, and a continuous health capital stock with diminishing-returns investment technology. \citet{Lim2020} analyzes the case of South Korea, showing that public provision of health insurance alters aggregate saving behavior in an OLG framework with endogenous health risk---a result that resonates with our model's mechanism, in which the scope of public health coverage affects both household savings decisions and the fiscal cost of the pension system through the contribution base.

%% --------------------------------------------------------------------
\subsection{Endogenous Health Investment and the Longevity Channel}
\label{subsec:lit_longevity}

A recent and rapidly developing strand of the literature studies health investment not merely as a household choice but as a macroeconomic transmission channel. \citet{Carrasco2026} provides a particularly relevant contribution for our framework. In his model, households' health investments endogenously determine longevity, which in turn affects aggregate savings, capital accumulation, and the effectiveness of macroeconomic policy. The key insight is that health investment creates a \textit{longevity channel}: when households invest more in health, they live longer, which strengthens the life-cycle savings motive, raises the aggregate capital stock, and alters the equilibrium interest rate. This channel operates in our model through the interaction of the health formation equation---where health investment $m_j$ augments the stock $h$ that governs both survival probabilities and labor productivity---and the capital market clearing condition, which equates domestic household savings to the aggregate capital stock plus government debt.

The relevance of the longevity channel for Latin America is amplified by two structural features. First, the premature health crossing point documented in Section~\ref{sec:intro}---the fact that health deterioration begins a decade earlier than in high-income countries---implies that the longevity channel operates asymmetrically across skill groups: low-skilled workers, who face earlier and steeper health depreciation, have shorter effective planning horizons and accumulate less precautionary savings, reinforcing the health-income poverty trap that our simulations identify. Second, the interaction between endogenous longevity and the PAYG pension system creates a fiscal amplification mechanism: healthier workers who live longer draw pension benefits for more periods, raising the system dependency ratio $N^R/N^W$ and the endogenous contribution rate $\tau^p$, which in turn increases the tax wedge on formal labor income. This feedback loop---from health investment to longevity to pension costs to the fiscal wedge---is the channel through which spending smarter on health can yield fiscal dividends beyond the direct reduction in out-of-pocket expenditures.

%% --------------------------------------------------------------------
\subsection{Informality, Formalization Incentives, and Social Security}
\label{subsec:lit_informality}

The empirical literature on formalization incentives in Latin America documents a robust negative relationship between social security contribution rates and formal employment, particularly at the lower end of the wage distribution. \citet{Kugler2017} estimate payroll tax elasticities in Colombia ranging between 0.4 and 1.0; similar magnitudes have been found in studies of Turkey \citep{AcarTansel2019}, Brazil, and Mexico. The perceived actuarial value of pension contributions---not just their cost---shapes formalization decisions: evidence from Colombia shows that an increase in the vesting period required to qualify for the minimum pension reduced formal employment by approximately 3 percentage points, illustrating that workers respond to the expected return on their contributions, not merely to the contribution rate itself.\footnote{This result is consistent with the ``social security wealth'' concept developed by \citet{WorldBankInformalizing2025}, which measures the actuarial present value of the benefit stream a worker expects to receive conditional on maintaining formal-sector contributions. Where social security wealth is low---because of long vesting periods, benefit opacity, or low institutional trust---the FTR rises and formalization incentives weaken.}

The non-contributory benefit channel operates in the opposite direction: policies that make health or income transfers available outside formal employment raise the opportunity cost of formalization and tend to shift labor toward the informal sector. \citet{BoschCampos2014} and \citet{AteridoEtAl2011} document significant informality effects from Mexico's Seguro Popular program, though results are not uniform---\citet{AzuaraMarinescu2013} find smaller or negligible effects, suggesting that institutional context and program design mediate the relationship. \citet{CamachoConoverHoyos2013} find analogous results for Colombia's subsidized health insurance expansion, with informality increasing by approximately 4 percentage points. \citet{CalderonMarinescu2012} present a more nuanced picture for Colombia, finding that regulations unifying health and pension payment systems partially offset the informality-inducing effects of coverage expansion. The general lesson, consistent with the analytical framework of this paper, is that the design of the health benefit package---in particular the degree to which it is tied to formal contributions versus financed through general taxation---is a first-order determinant of formalization equilibria.

The broader structural context is provided by \citet{PerryEtAl2007}, who argue that widespread informality reflects a mass opting out of state institutions---a ``blunt societal indictment of the quality of the State's services provision.'' \citet{WorldBankInformalizing2025} extend this line of analysis to its most comprehensive treatment to date, arguing that three dimensions must move together for formalization to improve: market competitiveness (the tax wedge must not exceed what firms can absorb), workers' valuation of contributory benefits (``social security wealth'' must be perceived as positive), and enforcement credibility (compliance must be monitored and sanctioned). Their reform roadmap emphasizes that LAC countries face a compressed timeline---the aging that took Europe 70 years will take Latin America approximately 25---leaving far less room for incremental adjustment. The report's conclusion that the ``crisis is not yet averted'' provides the institutional backdrop against which our model's policy counterfactuals should be read. \citet{BoschMelguizoPages2013} document how pension design interacts with labor market outcomes across the region, showing that systems with tighter benefit-contribution links generate stronger formalization incentives.

Our model contributes to this literature by calibrating the baseline to the formal-sector population and studying informality through a structured policy experiment: a formalization program that incorporates a fraction of the informal population into the contributory system, using observed income and health profiles from national household surveys. This design allows us to quantify the competing fiscal effects of a broader contribution base against higher per-capita health expenditure claims from newly formalized low-skilled workers, within the general-equilibrium discipline of endogenous factor prices and pension clearing.

%% --------------------------------------------------------------------
\subsection{Gender Heterogeneity in Lifecycle Fiscal Models}
\label{subsec:lit_gender}

The incorporation of gender into quantitative OLG models remains at an early stage. \citet{FehrKindermann2018} provide the most widely used pedagogical framework, constructing a two-person household in which the female member's labor force participation decision is modeled as a function of human capital accumulation, care responsibilities, and the relative returns to market work versus home production. Their model demonstrates the basic mechanism---female participation responds to policy instruments such as childcare subsidies, parental leave, and joint versus individual taxation---but operates in a substantially simpler economic environment than the one developed in this paper: health is not a state variable, there is no endogenous health investment or survival risk, no longevity channel, and no formalization margin.

The computational obstacle to combining household bargaining with rich health dynamics is severe. \citet{DeNardi2010} and subsequent work in the same line have extended models of medical expense risk to two-person households---singles and couples---but in a partial-equilibrium setting, with factor prices treated as exogenous. Embedding such a two-person household structure in a general-equilibrium framework, where wages and the interest rate respond endogenously to aggregate savings and labor supply, raises the computational burden substantially: the curse of dimensionality compounds with the fixed-point problem in factor prices. With the continuous health capital, age-dependent depreciation, endogenous survival, and endogenous factor prices that characterize our model, a structural household bargaining extension would be computationally challenging and, to our knowledge, has not been attempted in the published literature.

Our approach to gender is different. Rather than modeling the household decision-making process, we exploit the fact that sex-differentiated income profiles and survival probabilities already embed---in reduced form---the cumulative effects of labor market discrimination, career interruptions from child-bearing, gendered sorting into formal and informal employment, and the unequal distribution of care responsibilities. By expanding the agent-type space from two skill types $(\theta_L, \theta_H)$ to four $(\theta_L^M, \theta_H^M, \theta_L^F, \theta_H^F)$, calibrated with sex-specific earnings trajectories from household surveys and the survival differentials documented in national life tables, the model can decompose fiscal sustainability pressures by sex within the existing equilibrium framework. This calibration-based strategy lets the data carry the information that a structural bargaining model would generate endogenously, at a fraction of the computational cost and with direct empirical grounding. The structural household model remains a natural extension for future research, but it is not a prerequisite for the first-order fiscal decomposition that this paper's policy questions require.

\bigskip
%% --------------------------------------------------------------------

\section{Model}
\label{sec:model}

\bigskip

Section~\ref{sec:model} presents the theoretical framework that integrates these five strands. The model is a stochastic OLG economy with endogenous health investment, heterogeneous agents, and an explicit PAYG pension system; it extends the ITED baseline \citep{CortesEtAl2024} to incorporate general equilibrium factor pricing, a continuous health capital stock with age-dependent depreciation, and a contribution-rate rule that endogenously links demographic aging to the fiscal wedge on formal labor.
\subsection{Environment and Demographics}
\label{subsec:environment}

Time is discrete. The economy is populated by $J$ overlapping generations indexed by $j = 1, \ldots, J$, where $j = 1$ corresponds to age~20 and each model period spans five years. Households enter economic life at $j = 1$, retire exogenously at age $j_R$, and may live until $j = J$ (age~100).

\paragraph{Agent types.}
Each individual is indexed by sex $g \in \{m, f\}$ (male, female) and by a permanent skill type $\theta \in \{\theta_L, \theta_H\}$ (low-skilled, high-skilled), where $\theta_L$ and $\theta_H$ denote the log-productivity intercepts of the two skill groups (see equation~\eqref{eq:labor_productivity}). Sex and skill are both drawn at birth and remain fixed over the life cycle. Sex enters the model through two channels: a sex-specific age-efficiency profile $e^g_j$ that shapes earnings, and a sex-specific survival schedule $\psi^g_j$ that shapes longevity. Skill enters through the log-intercept $\theta$ and through a skill-specific health-productivity slope $\varrho(\theta)$. Let $\pi^{g,\theta}_1$ denote the birth share of agents of sex $g$ and skill $\theta$, with $\sum_{g,\theta} \pi^{g,\theta}_1 = 1$. The population grows at gross rate $1 + n_p$, so the absolute mass of newborns of type $(g,\theta)$ at calendar time $t$ is $\pi^{g,\theta}_1 \cdot N_t$.

\paragraph{Idiosyncratic risk.}
Households face two sources of individual uncertainty: (i)~mortality risk, depending on health, and (ii)~idiosyncratic labor-productivity risk, summarized by an AR(1) shock $\eta$. Health $h$ is endogenous: it evolves deterministically given current health and chosen medical spending. There is no aggregate uncertainty.

\paragraph{Individual state.}
The individual state of a household at the beginning of a period is
\begin{equation}
\label{eq:state}
s = (a,\, h,\, \eta),
\end{equation}
where $a \in \mathcal{A} = [0, \bar{a}]$ are beginning-of-period assets, $h \in \mathcal{H}_h = [0, \bar{h}]$ is health capital, and $\eta \in \mathcal{H}_\eta = \{\eta_1, \ldots, \eta_{N_\eta}\}$ is the idiosyncratic productivity component. Sex $g$, skill type $\theta$ and age $j$ are deterministic indices that label the value and policy functions.

\paragraph{Cohort weights.}
Let $\phi_j^{g}(\theta)$ denote the stationary mass of age-$j$ households of sex $g$ and skill $\theta$, normalized so that $\sum_g \sum_\theta \sum_{j=1}^{J} \phi_{j}^{g}(\theta) = 1$ with newborn weights $\phi_{1}^{g}(\theta) = \pi^{g,\theta}_1 / Z$, where
\[
Z \;=\; \sum_{g}\sum_{\theta}\sum_{j=1}^{J} \pi^{g,\theta}_1 \prod_{i=2}^{j} \Bigl[\bar\psi_i^{g}(\theta) \big/ (1+n_p)\Bigr]
\]
is the normalizing constant (with the convention that the empty product at $j=1$ equals one). Birth shares are sex- and skill-specific and need not coincide across types; calibration imposes $\pi^{m,\theta_H}_1 > \pi^{f,\theta_H}_1$, consistent with the data. The weights satisfy the forward recursion
\begin{equation}
\label{eq:cohort_weights}
\phi_{j+1}^{g}(\theta) \;=\; \frac{\bar{\psi}_{j+1}^{g}(\theta)}{1 + n_p}\, \phi_{j}^{g}(\theta), \qquad j = 1, \ldots, J-1,
\end{equation}
where the population-weighted average survival is
\begin{equation}
\label{eq:avg_survival}
\bar{\psi}_{j+1}^{g}(\theta)
\;=\;
\frac{1}{\phi_{j}^{g}(\theta)} \int \psi_{j+1}^{g}\bigl(h'(s; j, g, \theta)\bigr)\, d\mu_j^{g,\theta}(s),
\end{equation}
with $h'(s; j, g, \theta) = (1-\delta_j^h)\, h + H_j\bigl(m_j(s;\, g, \theta)\bigr)$ the endogenous next-period health and $\mu_{j}^{g,\theta}$ the unconditional age-$j$ population measure for sex $g$ and skill $\theta$ (so that $\int d\mu_{j}^{g,\theta} = \phi_{j}^{g}(\theta)$). Terminal-age survival is zero: $\psi_{J+1}^{g}(\cdot) \equiv 0$.

\paragraph{Within-period timing.}
The timing within each period is:
\begin{enumerate}
\item The household observes its current state $(a, h, \eta)$.
\item Working-age households choose $(c, \ell, a', m)$. Retired households choose $(c, a', m)$.
\item End-of-period assets are $a'$. Survival to next period is realized with probability $\psi_{j+1}^{g}(h')$, where $h' = (1-\delta_j^h) h + H_j(m)$ is determined by the chosen $m$. Assets of agents who die are voided: they neither enter aggregate capital next period nor flow to survivors. This rules out an accidental-bequest channel and isolates the life-cycle savings motive; the aggregate resource implication is tracked by the void term $\Lambda_{\mathrm{void}}$ introduced in equation~\eqref{eq:leakage_voided}.
\end{enumerate}


\subsection{Preferences}
\label{subsec:preferences}

Individuals rank streams of consumption, leisure, and health by expected discounted lifetime utility. The lifetime objective of a newborn household of sex $g$ type~$\theta$  is
\begin{equation}
\label{eq:lifetime_utility}
\E_1 \left[\sum_{j=1}^{J} \beta^{j-1} \left(\prod_{i=2}^{j} \psi_{i}^{g}(h_{i-1}')\right) u(c_j,\, \ell_j,\, h_j)\right],
\end{equation}
where $\beta \in (0,1)$, the survival probability $\psi_{i}^{g}(h_{i-1}')$ is the probability of being alive at age $i$ given that next-period health entering age $i$ is $h_{i-1}' = h_i$, and the expectation is taken over the random path $\{\eta_j\}_{j=1}^{J}$ of the idiosyncratic shock generated by~\eqref{eq:ar1}.
 
The period utility function is GHH-augmented with a health amenity:
\begin{equation}
\label{eq:utility}
u(c, \ell, h) \;\equiv\; \frac{\bigl(c + s(h) - v(\ell)\bigr)^{1-\gamma} - 1}{1 - \gamma},
\qquad
v(\ell) \equiv \Psi \frac{\ell^{1+\nu_\ell}}{1+\nu_\ell},
\end{equation}
\begin{equation}
    s(h) \equiv
\begin{cases}
\Xi \dfrac{h^{1-\xi}}{1-\xi} & \xi \neq 1,\\[6pt]
\Xi \log(h) & \xi = 1,
\end{cases}
\end{equation}

with $\Psi > 0$, $\Xi \geq 0$, $\nu_\ell > 0$ (the inverse Frisch elasticity, written with an $\ell$ subscript to avoid collision with the productivity function $\nu_j$), and $\gamma > 0$. Here $c \geq 0$ is consumption, $\ell \in [0, 1]$ is hours worked ($1-\ell$ is leisure), and $h$ is the contemporaneous stock of health capital that simultaneously affects survival, productivity, and the amenity value of consumption.

\begin{remark}[GHH composite with health amenity]
\label{rem:ghh_amenity}
Placing $s(h)$ inside the GHH composite makes the marginal utility of consumption depend on health, $u_{c,h} \neq 0$, so health investment and consumption are not separable at the FOC level. This is the channel through which $s(h)$ feeds into the medical-spending Euler equation~\eqref{eq:foc_m}; the standard wealth-effect-free property of GHH continues to hold in the intratemporal labor margin~\eqref{eq:cons_leisure}.
\end{remark}

\begin{remark}[Retirees]
\label{rem:retirees}
For retirees ($j \geq j_R$), $\ell_j = 0$ and the period utility reduces to $u(c, 0, h) = \frac{(c + s(h))^{1-\gamma} - 1}{1-\gamma}$.
\end{remark}


\subsection{Technology}
\label{subsec:technology}

A representative, competitive firm produces the single consumption good using capital $K$ and aggregate effective labor $L$ according to the Cobb-Douglas technology
\begin{equation}
\label{eq:production}
Y = A\, K^\alpha\, L^{1-\alpha}, \qquad \alpha \in (0,1),
\end{equation}
with $A > 0$ total factor productivity. Capital depreciates at rate $\delta \in (0,1)$. Because the production function has constant returns to scale, the firm's first-order conditions depend only on the ratio $K/L$ and so are invariant to the choice of units. We will write $K$ and $L$ in stationary per-cohort-mass units throughout (see notation block in Section~\ref{subsec:markets}); this convention is consistent with the firm's optimization.
 
\paragraph{Firm's problem.}
The firm takes factor prices $(r, w)$ as given and maximizes static profits:
\begin{equation}
\label{eq:firm_problem}
\max_{K, L} \bigl\{A K^\alpha L^{1-\alpha} - (r+\delta)\,K - w\,L\bigr\}.
\end{equation}
The first-order conditions yield
\begin{align}
r &= \alpha A \left(\frac{K}{L}\right)^{\alpha-1} - \delta, \label{eq:foc_r} \\
w &= (1-\alpha)\, A \left(\frac{K}{L}\right)^{\alpha}. \label{eq:foc_w}
\end{align}

\paragraph{Factor prices.}
Both $r$ and $w$ are determined endogenously in equilibrium: given aggregate capital $K$ and aggregate effective labor $L$, the capital--labor ratio $K/L$ pins down both prices simultaneously via~\eqref{eq:foc_r}--\eqref{eq:foc_w}. The interest rate is the object that adjusts to clear the capital market.


\subsection{Exogenous Processes}
\label{subsec:exogenous}

\subsubsection{Mortality}
\label{subsubsec:mortality}

Conditional survival probabilities depend on age, sex, and end-of-period health. The probability of surviving from age $j$ to age $j+1$, realized at the end of age $j$, is $\psi_{j+1}^{g}(h')$, where $h'$ is the next-period health stock determined by the health-formation rule~\eqref{eq:health_formation}. Terminal-age survival is zero, $\psi_{J+1}^{g}(\cdot) \equiv 0$, so death by age $j = J$ is certain.

\subsubsection{Labor Productivity}
\label{subsubsec:productivity}

Effective labor productivity of a working-age household ($j < j_R$) is
\begin{equation}
\label{eq:labor_productivity}
\nu_j(h, \eta;\, g,\theta) \;=\; e_{j}^{g} \exp\bigl(\theta + \eta + \varrho(\theta)\, h\bigr),
\end{equation}
where $e_{j}^{g} > 0$ is a deterministic sex-dependent age-efficiency profile, $\theta \in \{\theta_L, \theta_H\}$ is the log-productivity intercept of the skill group, and $\varrho(\theta) \geq 0$ is the (skill-specific) elasticity of productivity with respect to health capital. Productivity is zero for retirees ($j \geq j_R$).

The idiosyncratic component $\eta$ follows the AR(1) process
\begin{equation}
\label{eq:ar1}
\eta_{j+1} = \rho\, \eta_j + \varepsilon_{j+1}, \qquad \varepsilon_{j+1} \sim \mathcal{N}(0, \sigma_\varepsilon^2),
\end{equation}
with $|\rho| < 1$. 

\subsubsection{Health Formation and Health Investment}
\label{subsubsec:health}

Health evolves deterministically given the chosen medical spending:
\begin{equation}
\label{eq:health_formation}
h' = \min\bigl\{(1 - \delta_j^h)\, h + H_j(m),\; \bar h\bigr\},
\end{equation}
where the depreciation rate $\delta_j^h \in (0, 1)$ is age-dependent and $H_j: \mathbb{R}_+ \to \mathbb{R}_+$ satisfies $H_j(0) = 0$, $H_j' > 0$, $H_j'' < 0$ (diminishing returns to medical spending). A specific functional form (e.g., $H_j(m) = \xi_j m^{\zeta}$ with $\zeta \in (0,1)$) must be committed to before computation.


\subsection{Government}
\label{subsec:government}

The government runs two separate balance sheets: a general fiscal block (taxes, public goods, and one-period debt) and a self-financed pay-as-you-go (PAYG) pension block. Pension contributions $\tau^p$ enter the general budget only indirectly, through the labor-tax wedge faced by workers; they are not a revenue item for the general government.

\subsubsection{Fiscal Policy}
\label{subsubsec:fiscal}
The government taxes consumption at rate $\tau^c$, labor income at rate $\tau^\omega$, capital income at rate $\tau^k$, and medical spending at rate $\tau^m$. It purchases public goods $G = g_y Y$ and may issue one-period debt $B$. All aggregate quantities ($C$, $K$, $L$, $M$, $Y$, $G$, $B$) are in stationary per-cohort-mass units (see notation block in Section~\ref{subsec:markets}); $C$, $L$, $M$, $Y$, $G$ are per-period flows and $K$, $B$ are stocks. The flow budget is:
\begin{equation}
\label{eq:gov_budget}
\tau^c C + \tau^m M + \tau^\omega w L + \tau^k r K + B' \;=\; G + (1+r)B,
\end{equation}
where $C$ is aggregate consumption, $K$ is the aggregate capital stock, $M$ is aggregate health investments, and $B'$ denotes next-period government debt.

\paragraph{Fiscal closure.}
$(\tau^c, \tau^\omega, \tau^k, \tau^m, g_y)$ are exogenous policy parameters. In the stationary equilibrium $B' = B$, so debt service $r B$ equals the primary surplus:
\begin{equation}
\label{eq:gov_stationary}
r\, B \;=\; \tau^c\, C + \tau^\omega\, w\, L + \tau^k\, r\, K + \tau^m\, M - G.
\end{equation}
$B$ is the residual that closes the budget.

\subsubsection{Pension System}
\label{subsubsec:pension}
A PAYG pension is financed by a payroll tax $\tau^p$ on workers' labor income. The system balances each period:
\begin{equation}
\label{eq:pension_balance}
\tau^p\, w\, L \;=\; \overline{pen}\, N^R,
\end{equation}
where $N^R = \sum_g \sum_\theta \sum_{j=j_R}^{J} \phi_j^{g}(\theta)$ is the mass of retirees. The pension benefit is the replacement-rate parameter $\kappa$ times average per-worker labor income:
\begin{equation}
\label{eq:pension_benefit}
\overline{pen} \;=\; \kappa \cdot \frac{w L}{N^W},
\end{equation}
where $N^W = \sum_g \sum_\theta \sum_{j=1}^{j_R - 1} \phi_j^{g}(\theta)$ is the mass of working-age individuals.

\paragraph{Pension closure.}
The replacement rate $\kappa$ is an exogenous policy parameter. The contribution rate $\tau^p$ adjusts endogenously to balance the pension system each period. Substituting~\eqref{eq:pension_benefit} into~\eqref{eq:pension_balance} yields
\begin{equation}
\label{eq:pension_contribution}
\tau^p \;=\; \kappa\, \frac{N^R}{N^W},
\end{equation}
so the contribution rate equals the replacement rate scaled by the system dependency ratio $N^R / N^W$. This relationship pins $\tau^p$ down once the stationary distribution is known.


\subsection{Markets}
\label{subsec:markets}
\par Throughout this section we work in \emph{stationary per-cohort-mass} units. Aggregates are integrals of individual quantities against the stationary population measure $\mu_j^{g,\theta}$, with $\sum_g \sum_\theta \sum_j \phi_j^{g}(\theta) = 1$. Population growth $n_p$ is already absorbed into the cohort-weight recursion~\eqref{eq:cohort_weights}; the aggregates defined below therefore do \emph{not} grow with the actual population. A symbol such as $K$ in the Markets and Equilibrium sections refers to the stationary capital stock per unit of cohort mass, related to the absolute capital stock by $K^{\mathrm{abs}}_t = K \cdot N_t$, where $N_t$ is total population at $t$ growing at gross rate $1+n_p$. We distinguish stocks from per-period flows. \emph{Stocks} (units: goods, or units of agent-mass): $K$, $A_{\mathrm{dom}}$, $B$, $N^W$, $N^R$. \emph{Per-period flows} (units: goods per period, or efficiency units of labor per period): $Y$, $C$, $M$, $G$, $(n_p+\delta) K$, $L$, $\Lambda_{\mathrm{void}}$, $\overline{pen}$. Gross investment in capital is $(n_p+\delta) K$ per cohort-mass per period: $\delta K$ replaces depreciation and $n_p K$ provides for the next cohort that enters at gross rate $1+n_p$. The goods-market identity~\eqref{eq:goods_market} is a flow identity: every term has units of goods per period.

\par Asset markets are incomplete. Households can save in a single non-contingent asset $a$ and are not allowed to borrow, so
\begin{equation}
\label{eq:borrowing_constraint}
a' \;\geq\; 0.
\end{equation}
Both factor prices $(r, w)$ are determined endogenously by equations~\eqref{eq:foc_r}--\eqref{eq:foc_w}. Markets clear when the following conditions hold:

\paragraph{Labor market.}
\begin{equation}
\label{eq:labor_market}
L \;=\; \sum_g\, \sum_\theta\, \sum_{j=1}^{j_R - 1}\, \int \nu_j(h, \eta;\, g, \theta)\, \ell_j(s;\, g, \theta)\, d\mu_{j}^{g,\theta}(s).
\end{equation}
\paragraph{Capital market.}
Domestic household savings finance the entire capital stock and government debt:
\begin{equation}
\label{eq:capital_market}
K \;=\; A_{\mathrm{dom}} - B,
\qquad
A_{\mathrm{dom}} \;=\; \sum_g\,\sum_\theta\, \sum_{j=1}^{J}\, \int a\, d\mu_j^{g,\theta}(s).
\end{equation}
This condition, together with the firm's FOC~\eqref{eq:foc_r}, determines the equilibrium interest rate.

\paragraph{Goods market.}
\begin{equation}
\label{eq:goods_market}
Y \;=\; C + M + (n_p + \delta)\, K + G + \Lambda_{\mathrm{void}},
\end{equation}
where every term is a flow (goods per period in stationary per-cohort-mass units), $(n_p + \delta) K$ is gross investment as discussed in the notation block, and the voided-asset flow is
\begin{equation}
\label{eq:leakage_voided}
\Lambda_{\mathrm{void}}
\;\equiv\;
\sum_g\, \sum_\theta\, \sum_{j=1}^{J}\, \bigl(1 - \bar\psi_{j+1}^{g}(\theta)\bigr)
\int a_j'(s;\, g, \theta)\, d\mu_{j}^{g,\theta}(s).
\end{equation}
\paragraph{Government bond market.}
By construction the government issues exactly $B$ bonds; households' aggregate asset holdings finance both $K$ and $B$:
\begin{equation}
A_{\mathrm{dom}} = K + B,
\end{equation}
which is~\eqref{eq:capital_market} rearranged.

\subsection{Optimization Problems}
\label{subsec:optimization}

\subsubsection{Net Labor Income}
\label{subsubsec:net_income}

Define net labor income at age $j$ as
\begin{equation}
\label{eq:net_labor_income}
y_j(h, \eta, \ell;\, g, \theta) \;=\; w\, \nu_j(h, \eta;\, g, \theta)\, \ell\, \bigl(1 - \tau^\omega - \tau^p\bigr),
\end{equation}
where $\tau^\omega$ is the labor income tax and $\tau^p$ is the pension contribution rate. All workers in the baseline model are formal; informality is studied as a calibration-based counterfactual in Section~\ref{sec:formalization}.

\subsubsection{Workers' Problem}
\label{subsubsec:worker}

Individuals observe their productivity shock and then choose consumption, savings, labor supply, and health investment. Given labor supply $\ell$, define household resources at age $j$ as
\begin{equation}
\label{eq:resources_worker}
x_j(a, h, \eta, \ell;\, g, \theta) \;=\; \bigl(1 + r(1 - \tau^k)\bigr) a + y_j(h, \eta, \ell;\, g, \theta).
\end{equation}
The worker's Bellman equation is
\begin{equation}
\label{eq:worker_bellman}
V_j(a, h, \eta;\, g,\theta)
\;=\;
\max_{\substack{c \geq 0,\, a' \geq 0 \\ \ell \in [0,1],\, m \geq 0}}
\Bigl\{
u(c, \ell, h)
+ \beta\, \psi_{j+1}^{g}(h')\,
\E\bigl[V_{j+1}(a', h', \eta';\, g, \theta) \mid \eta\bigr]
\Bigr\}
\end{equation}
subject to
\begin{align}
(1 + \tau^c)\, c + (1 + \tau^m)\, m + a' &= x_j(a, h, \eta, \ell;\, g, \theta), \label{eq:worker_budget}\\
h' &= \min\bigl\{(1 - \delta_j^h)\, h + H_j(m),\, \bar h\bigr\}, \label{eq:worker_health}\\
a' &\geq 0,\quad m \geq 0,\quad \ell \in [0,1]. \label{eq:worker_borrow}
\end{align}
Let the optimal savings policy be $a_j'(a, h, \eta;\, g, \theta)$ and analogously for $c_j, \ell_j, m_j$.

\subsubsection{Retiree's Problem}
\label{subsubsec:retiree}
For retirees ($j \geq j_R$) labor income is zero, so resources are
\begin{equation}
\label{eq:resources_retiree}
x_j(a;\, g, \theta) \;=\; \bigl(1 + r(1-\tau^k)\bigr) a + \overline{pen}.
\end{equation}
The productivity shock $\eta$ does not enter retirees' resources but is retained as a state to keep the state space dimension uniform across ages. The Bellman equation is
\begin{equation}
\label{eq:retiree_bellman}
V_j(a, h, \eta;\, g, \theta)
\;=\;
\max_{c \geq 0,\, a' \geq 0,\, m \geq 0}
\Bigl\{
u(c, 0, h)
+ \beta\, \psi_{j+1}^{g}(h')\,
\E\bigl[V_{j+1}(a', h', \eta';\, g, \theta) \mid \eta\bigr]
\Bigr\}
\end{equation}
subject to
\begin{align}
(1 + \tau^c)\, c + (1 + \tau^m)\, m + a' &= x_j(a;\, g, \theta), \label{eq:retiree_budget}\\
h' &= \min\bigl\{(1 - \delta_j^h)\, h + H_j(m),\, \bar h\bigr\}, \label{eq:retiree_health}\\
a' &\geq 0,\quad m \geq 0. \label{eq:retiree_borrow}
\end{align}
At the terminal age $j = J$, $\psi_{J+1}^{g}(h') = 0$, so the continuation value vanishes and the optimum is $a'_J = 0$, $m_J = 0$ (no point investing in health when survival probability is zero).

\subsubsection{First-Order Conditions}
\label{subsubsec:focs}

Let $\lambda_j$ denote the multiplier on the budget constraint. The interior FOCs are:
\begin{align}
\text{Consumption:} \quad & u_c(c, \ell, h) \;=\; \lambda_j (1 + \tau^c), \label{eq:foc_c}\\
\text{Labor:} \quad & -u_\ell(c, \ell, h) \;=\; \lambda_j\, w\, \nu_j(h, \eta;\, g, \theta)\,(1-\tau^\omega-\tau^p), \label{eq:foc_l}\\
\text{Savings:} \quad & \lambda_j \;=\; \beta\, \psi_{j+1}^{g}(h')\, \E\bigl[V_{a, j+1}(a', h', \eta';\, g, \theta) \mid \eta\bigr], \label{eq:foc_a}\\
\text{Medical spending:} \quad & \lambda_j (1 + \tau^m)
\;=\; \beta\, H_j'(m)\,\Bigl\{
\bigl(\psi_{j+1}^{g}\bigr)'(h')\, \E[V_{j+1}(a', h', \eta';\, g, \theta) \mid \eta] \notag\\
&\qquad\qquad\qquad\qquad +\, \psi_{j+1}^{g}(h')\, \E[V_{h, j+1}(a', h', \eta';\, g, \theta) \mid \eta]
\Bigr\}. \label{eq:foc_m}
\end{align}

\paragraph{Envelope conditions.} For the worker's problem ($j < j_R$),
\begin{align}
V_{a, j}(a, h, \eta;\, g, \theta) &\;=\; \lambda_j \cdot \bigl(1 + r(1-\tau^k)\bigr), \label{eq:env_a}\\
V_{h, j}(a, h, \eta;\, g, \theta) &\;=\; u_h(c, \ell, h) + \lambda_j\, w\, \nu_j(h, \eta;\, g, \theta)\, \varrho(\theta)\, \ell\, (1-\tau^\omega-\tau^p) \notag\\
&\qquad +\, \beta\, (1-\delta_j^h)\, \bigl(\psi_{j+1}^{g}\bigr)'(h')\, \E[V_{j+1}(a', h', \eta';\, g, \theta) \mid \eta] \notag\\
&\qquad +\, \beta\, (1-\delta_j^h)\, \psi_{j+1}^{g}(h')\, \E[V_{h, j+1}(a', h', \eta';\, g, \theta) \mid \eta]. \label{eq:env_h}
\end{align}
For the retiree's problem ($j \geq j_R$), $\ell = 0$ and the labor-income term in $V_h$ vanishes:
\begin{align}
V_{a, j}^{R}(a, h, \eta;\, g, \theta) &\;=\; \lambda_j \cdot \bigl(1 + r(1-\tau^k)\bigr), \label{eq:env_a_ret}\\
V_{h, j}^{R}(a, h, \eta;\, g, \theta) &\;=\; u_h(c, 0, h) \notag\\
&\qquad +\, \beta\, (1-\delta_j^h)\, \bigl(\psi_{j+1}^{g}\bigr)'(h')\, \E[V_{j+1}(a', h', \eta';\, g, \theta) \mid \eta] \notag\\
&\qquad +\, \beta\, (1-\delta_j^h)\, \psi_{j+1}^{g}(h')\, \E[V_{h, j+1}(a', h', \eta';\, g, \theta) \mid \eta]. \label{eq:env_h_ret}
\end{align}
At the terminal age $j = J$, $\psi^{g}_{J+1} \equiv 0$ kills the continuation terms; the envelopes are still well defined since $a_J' = m_J = 0$.

\paragraph{Euler equation.} Combining \eqref{eq:foc_a} with \eqref{eq:env_a} (one period ahead) and \eqref{eq:foc_c}:
\begin{equation}
\label{eq:euler}
u_c(c, \ell, h) \;=\; \beta\, \bigl(1 + r(1-\tau^k)\bigr)\, \psi_{j+1}^{g}(h')\, \E\bigl[u_c(c', \ell', h') \mid \eta\bigr].
\end{equation}
 
\paragraph{Consumption--leisure condition.} Combining \eqref{eq:foc_c} and \eqref{eq:foc_l}:
\begin{equation}
\label{eq:cons_leisure}
\frac{-u_\ell(c, \ell, h)}{u_c(c, \ell, h)}
\;=\;
\frac{w\, \nu_j(h, \eta; g, \theta)(1-\tau^\omega-\tau^p)}{1 + \tau^c}.
\end{equation}
 
\paragraph{Kuhn--Tucker for the borrowing constraint.}
The borrowing constraint $a' \geq 0$ generates a non-negative multiplier and turns the savings FOC~\eqref{eq:foc_a} into the weak inequality
\begin{equation}
\label{eq:kt_borrowing}
\lambda_j \;\geq\; \beta\, \psi_{j+1}^{g}(h')\, \E\bigl[V_{a, j+1}(a', h', \eta';\, g, \theta) \mid \eta\bigr], \qquad a' \geq 0,
\end{equation}
with complementary slackness $\bigl(\lambda_j - \beta\, \psi_{j+1}^{g}(h')\, \E[V_{a, j+1}(a', h', \eta';\, g, \theta) \mid \eta]\bigr)\, a' = 0$.

\paragraph{Kuhn--Tucker for medical spending.}
The non-negativity constraint $m \geq 0$ turns the medical-spending FOC~\eqref{eq:foc_m} into the weak inequality
\begin{equation}
\label{eq:kt_medical}
\lambda_j (1 + \tau^m)
\;\geq\;
\beta\, H_j'(m)\,\Bigl\{
\bigl(\psi_{j+1}^{g}\bigr)'(h')\, \E[V_{j+1}(\cdot) \mid \eta]
+ \psi_{j+1}^{g}(h')\, \E[V_{h, j+1}(\cdot) \mid \eta]
\Bigr\}, \qquad m \geq 0,
\end{equation}
with complementary slackness $\bigl(\lambda_j (1+\tau^m) - \beta H_j'(m)\{\cdot\}\bigr)\, m = 0$. The health cap $h' \leq \bar h$ in~\eqref{eq:health_formation} introduces a second non-negative multiplier $\mu_{\bar h} \geq 0$, with complementary slackness
\begin{equation}
\label{eq:kt_healthcap}
\mu_{\bar h}\,(\bar h - h') \;=\; 0, \qquad h' \leq \bar h, \qquad \mu_{\bar h} \geq 0,
\end{equation}
and $\mu_{\bar h}$ enters the medical-spending FOC as a subtracted term $\mu_{\bar h}\, H_j'(m)$ on the right-hand side of~\eqref{eq:kt_medical}. At an interior optimum with $m > 0$ and $h' < \bar h$, both inequalities hold with equality and~\eqref{eq:foc_m} is recovered.

\subsection{Stationary Competitive Equilibrium}
\label{subsec:equilibrium}

Let $\mu_j^{g,\theta}$ denote the time-invariant population measure of age-$j$ households of sex $g$ and skill $\theta$ over the state space $\mathcal{S} = \mathcal{A} \times \mathcal{H}_h \times \mathcal{H}_\eta$, with $\mathcal{A} = [0, \bar a]$, $\mathcal{H}_h = [0, \bar h]$, and $\mathcal{H}_\eta = \{\eta_1, \ldots, \eta_{N_\eta}\}$. The convention is that $\mu_{j}^{g,\theta}$ is \textbf{unconditional}: $\int d\mu_j^{g,\theta} = \phi_j^{g}(\theta)$, so aggregates are integrals over this measure without an outer cohort weight. Let $\Pi^\eta(\eta, \eta')$ denote the transition kernel of the AR(1) shock~\eqref{eq:ar1}, discretized to the finite grid $\mathcal{H}_\eta$ via Tauchen or Rouwenhorst (details deferred to the Calibration section).

\subsubsection{Aggregate Objects}

\begin{align}
C &\;=\; \sum_g\,\sum_\theta\, \sum_{j=1}^{J}\, \int c_j(s;\, g, \theta)\, d\mu_j^{g,\theta}(s), \label{eq:agg_C}\\
M &\;=\; \sum_g\,\sum_\theta\, \sum_{j=1}^{J}\, \int m_j(s;\, g, \theta)\, d\mu_j^{g,\theta}(s), \label{eq:agg_M}\\
L &\;=\; \sum_g\,\sum_\theta\, \sum_{j=1}^{j_R - 1}\, \int \nu_j(h, \eta;\, g, \theta)\, \ell_j(s;\, g, \theta)\, d\mu_j^{g,\theta}(s), \label{eq:agg_L}\\
A_{\mathrm{dom}} &\;=\; \sum_g\,\sum_\theta\, \sum_{j=1}^{J}\, \int a\, d\mu_j^{g,\theta}(s), \label{eq:agg_A}\\
N^W &\;=\; \sum_g\,\sum_\theta\, \sum_{j=1}^{j_R - 1}\, \phi_{j}^{g}(\theta), \qquad N^R \;=\; \sum_g\,\sum_\theta\, \sum_{j=j_R}^{J}\, \phi_{j}^{g}(\theta), \label{eq:agg_NWR}\\
\Lambda_{\mathrm{void}} &\;=\; \sum_g\,\sum_\theta\, \sum_{j=1}^{J}\, \bigl(1 - \bar\psi_{j+1}^{g}(\theta)\bigr) \int a_j'(s;\, g, \theta)\, d\mu_j^{g,\theta}(s). \label{eq:agg_Lvoided}
\end{align}

\subsubsection{Law of Motion for the Distribution}

For each measurable set $\mathcal{B} \subseteq \mathcal{S}$, each age $j = 1, \ldots, J-1$, each sex $g$, and each skill $\theta$:
\begin{equation}
\label{eq:lom}
\mu_{j+1}^{g,\theta}(\mathcal{B})
\;=\;
\frac{1}{1 + n_p}
\int
\psi_{j+1}^{g}\bigl(h'(s; j, g, \theta)\bigr)\,
\Bigl[\sum_{\eta' \in \mathcal{H}_\eta}
\Pi^\eta(\eta, \eta')\,
\mathbf{1}\bigl[\bigl(a_j'(s;\, g, \theta),\, h'(s; j, g, \theta),\, \eta'\bigr) \in \mathcal{B}\bigr]
\Bigr]
\, d\mu_j^{g,\theta}(s).
\end{equation}
The $1/(1+n_p)$ factor expresses the unconditional next-period measure in per-newborn-cohort units, consistent with the cohort-weight recursion~\eqref{eq:cohort_weights}. The recursion runs for $j = 1, \ldots, J-1$; no measure $\mu^{g,\theta}_{J+1}$ is defined because $\psi^{g}_{J+1} \equiv 0$ implies all age-$J$ agents die at the period boundary. The initial condition for newborns is
\begin{equation}
\label{eq:initial_dist}
\mu_1^{g,\theta}(\mathcal{B})
\;=\;
\phi_{1}^{g}(\theta) \cdot
\sum_{\eta_1 \in \mathcal{H}_\eta} \pi^\eta_{\mathrm{erg}}(\eta_1)\,
\mathbf{1}[(0, h_0, \eta_1) \in \mathcal{B}],
\end{equation}
where $\pi^\eta_{\mathrm{erg}}$ is the ergodic distribution of $\Pi^\eta$ and $h_0 \in (0, \bar h]$ is initial health (a parameter, e.g.\ $h_0 = \bar h$).

\subsubsection{Definition}
 
\begin{definition}[Stationary Recursive Competitive Equilibrium]
\label{def:equilibrium}
Given exogenous policy $(\tau^c, \tau^\omega, \tau^k, \tau^m, g_y, \kappa)$, birth shares $\{\pi^{g,\theta}_1\}_{g,\theta}$, population growth $n_p$, initial health endowment $h_0 \in (0, \bar h]$, and parameters of preferences, technology, mortality, productivity, and health formation, a \textbf{stationary recursive competitive equilibrium} consists of
\begin{enumerate}
\item value functions $\{V_j(\cdot;\, g, \theta)\}_{j=1,\ldots,J;\, g,\theta}$ and policy functions $\{c_j, \ell_j, m_j, a_j'\}(\cdot;\, g, \theta)$ defined on $\mathcal{S}$;
\item factor prices $(r, w)$;
\item endogenous fiscal/pension variables $(B, \tau^p, \overline{pen})$;
\item aggregates $(K, L, C, M, A_{\mathrm{dom}}, N^W, N^R, \Lambda_{\mathrm{void}})$ in stationary per-cohort-mass units (see notation block in Section~\ref{subsec:markets});
\item cohort weights $\{\phi_j^{g}(\theta)\}_{j, g, \theta}$ and population-average survivals $\{\bar\psi_{j+1}^{g}(\theta)\}_{j, g, \theta}$;
\item stationary distributions $\{\mu_j^{g,\theta}\}_{j, g, \theta}$;
\end{enumerate}
such that:
\begin{enumerate}
\item \textbf{Worker optimality.} For each $(g, \theta)$ and $j < j_R$, $V_j$ solves the worker's Bellman~\eqref{eq:worker_bellman}--\eqref{eq:worker_borrow}, with policy functions $(c_j, \ell_j, m_j, a_j')$.
\item \textbf{Retiree optimality.} For each $(g, \theta)$ and $j \geq j_R$, $V_j$ solves the retiree's Bellman~\eqref{eq:retiree_bellman}--\eqref{eq:retiree_borrow}, with policy functions $(c_j, m_j, a_j')$ and $\ell_j \equiv 0$.
\item \textbf{Firm optimality.} $(r, w)$ satisfy~\eqref{eq:foc_r}--\eqref{eq:foc_w} given $(K, L)$.
\item \textbf{Government budget.} $B$ satisfies~\eqref{eq:gov_stationary}.
\item \textbf{Pension balance.} $\tau^p$ satisfies~\eqref{eq:pension_contribution} and $\overline{pen}$ satisfies~\eqref{eq:pension_benefit}.
\item \textbf{Market clearing.} Labor~\eqref{eq:labor_market}, capital~\eqref{eq:capital_market}, and goods~\eqref{eq:goods_market}.
\item \textbf{Cohort weights.} $\phi_{j}^{g}(\theta)$ satisfy~\eqref{eq:cohort_weights} with $\bar\psi_{j+1}^{g}(\theta)$ given by~\eqref{eq:avg_survival} and birth shares $\phi_1^{g}(\theta) \propto \pi^{g,\theta}_1$, normalized so that $\sum_g \sum_\theta \sum_j \phi_j^{g}(\theta) = 1$.
\item \textbf{Stationarity of the distribution.} For each $(g, \theta, j)$, $\mu_j^{g,\theta}$ satisfies the law of motion~\eqref{eq:lom} with the newborn condition~\eqref{eq:initial_dist}.
\item \textbf{Aggregate consistency.} $(C, M, L, A_{\mathrm{dom}}, N^W, N^R, \Lambda_{\mathrm{void}})$ are computed from~\eqref{eq:agg_C}--\eqref{eq:agg_Lvoided}, all in stationary per-cohort-mass units.
\item \textbf{Demographic--distributional fixed point.} Conditions 7--8 are mutually dependent: the population-average survivals $\bar\psi^{g}_{j+1}(\theta)$ are functionals of the stationary distribution $\mu^{g,\theta}_j$ through~\eqref{eq:avg_survival}, while $\mu^{g,\theta}_j$ itself depends on $\bar\psi^{g}_{j+1}(\theta)$ through the $1/(1+n_p)$ scaling in~\eqref{eq:lom} that anchors the cohort weights. Equilibrium requires $\bigl(\{\phi^{g}_j(\theta)\}, \{\bar\psi^{g}_{j+1}(\theta)\}, \{\mu^{g,\theta}_j\}\bigr)$ to be a joint fixed point of~\eqref{eq:cohort_weights},~\eqref{eq:avg_survival}, and~\eqref{eq:lom}, evaluated at the equilibrium policy functions.
\end{enumerate}
\end{definition}

\subsection{Existence and Computation}
\label{subsec:existence}

The equilibrium of Definition~\ref{def:equilibrium} is the fixed point of three nested maps. The inner map is the household problem, solved by backward induction on the finite-horizon Bellman equations~\eqref{eq:worker_bellman} and~\eqref{eq:retiree_bellman} together with the savings and medical-spending first-order conditions~\eqref{eq:foc_a}--\eqref{eq:foc_m}. The middle map iterates the population measure~\eqref{eq:lom} jointly with the cohort weights~\eqref{eq:cohort_weights} and the health-dependent survival schedule~\eqref{eq:avg_survival} to their mutual fixed point---the demographic--distributional fixed point of the definition. The outer map adjusts the capital--labor ratio until the factor markets~\eqref{eq:foc_r}--\eqref{eq:foc_w} and the capital-market condition~\eqref{eq:capital_market} clear, with the contribution rate and the stationary debt level moving through~\eqref{eq:pension_contribution} and~\eqref{eq:gov_stationary}. The inner problem is solved by backward induction with constrained maximization over $(a',m)$---labour supply concentrated out through the intratemporal condition~\eqref{eq:cons_leisure}---on grids of $N_a=100$ asset, $N_h=15$ health, and $N_\eta=7$ productivity nodes for each of the $J=17$ ages and four types. We do not use the endogenous grid method: because medical spending enters the continuation value non-separably and through its \emph{level} (the survival derivative in~\eqref{eq:foc_m} multiplies $\E[V_{j+1}]$, not only its marginal value), and because next-period health is itself a control, the end-of-period state is not exogenous and the Euler equation cannot be inverted in closed form. The model admits a well-behaved stationary equilibrium of this form, and the policy functions deliver the lifecycle structure one expects of a life-cycle economy: consumption and assets are hump-shaped over the working life, medical spending rises through middle age and then flattens as the survival return to health investment falls, and the health stock deteriorates faster for low-skill types through the productivity penalty $\varrho(\theta)$. Assets of decedents are voided rather than bequeathed or annuitized (the leakage $\Lambda_{\mathrm{void}}$ in~\eqref{eq:goods_market}), which isolates the life-cycle saving motive; the resource leakage is a few percent of output, and disciplining it through an accidental-bequest or annuity closure is a robustness check we leave for the next stage. We display these profiles at the calibrated baseline in Section~\ref{subsec:calibrated_baseline}; their qualitative shape does not depend on the calibration.

Two features of the closure organize everything that follows. First, the PAYG block ties the contribution rate to the system dependency ratio, $\tau^p = \kappa\,N^R/N^W$, so any force that moves the ratio of retirees to workers moves the labor-tax wedge one for one. Second, because survival is endogenous to medical spending, health investment feeds forward into the dependency ratio itself: longer lives raise $N^R/N^W$, and hence $\tau^p$. The interaction of these two links---demographic pressure transmitted to the fiscal wedge, and health-driven longevity feeding back into that same wedge---is the \emph{fiscal amplification loop}. It is the mechanism the experiments of Section~\ref{sec:experiments} quantify, and the reason the net fiscal effect of better population health is ambiguous in sign.


%% ====================================================================
%% SECTION 4 — CALIBRATION (Mexican-data)
%% ====================================================================
\section{Calibration}
\label{sec:calibration}

This section is the empirical core of the paper. We calibrate the model to M\'exico in two steps. The first step fixes the model's exogenous processes and the externally assigned parameters: demographics, mortality, the sex-by-age labour-efficiency profile and skill premium estimated from the ENOE 2018--2024 panel, the productivity process, the health-depreciation schedule, and the fiscal block. The second step estimates the six remaining preference and health-technology parameters---which have no direct empirical counterpart---by simulated method of moments (SMM). The labour primitives are the team's own estimates from Mexican micro-data; the survival schedule, the health-depreciation schedule, the productivity process, and the medical-spending and value-of-statistical-life (VSL) moment targets are anchored to external data and the published literature for an upper-middle-income economy. The roadmap in Section~\ref{subsec:roadmap} states which inputs a forthcoming refinement will replace with additional Mexican micro-data; the present calibration is the basis for every number in this paper. To keep the economics in view, the SMM mechanics, the identification analysis, the convergence record, and the solver diagnostics are collected in Appendix~\ref{app:calibration}, and the three-country comparative landscape---which motivates the project but does not discipline the Mexican calibration---is collected in Appendix~\ref{app:supplementary}.

\subsection{Units and Normalization}
\label{subsec:units}

The model is dimensionless until two normalizations fix its scale, and every quantitative statement in the paper depends on stating them once and applying them consistently. \emph{Time.} A model period is five years; $j=1$ is age 20 and retirement is at $j_R=10$ (age 65). Output, consumption, medical spending, and the wage bill are five-year flows; capital and government debt are stocks. Ratios of a stock to a flow are therefore reported at annual frequency by multiplying the model (five-year) ratio by five: the annual capital--output ratio is $5\,(K/Y)$ and the annual debt--output ratio is $5\,(B/Y)$. Rates are annualized geometrically, $r^{\mathrm{ann}}=(1+r)^{1/5}-1$ for the five-year return $r$. \emph{Goods.} One unit of the consumption good equals USD~$35{,}000$, the level of per-capita household final consumption in M\'exico over a five-year period ($\approx$ USD~$7{,}000$ per year, 2021 current USD). This single factor, $\text{usd\_per\_unit\_c}=35{,}000$, converts every model-unit quantity---medical spending, the pension flow, the VSL---into dollars. \emph{Labour.} Effective labour is measured in efficiency units normalized so that the age-efficiency of a male aged 20--24 equals one; the wage $w$ is the price of one such efficiency unit. \emph{Population.} All aggregates are in stationary per-cohort-mass units (Section~\ref{subsec:markets}); levels therefore have no independent meaning across calibrations and only ratios, rates, and within-calibration comparisons are interpreted.

\subsection{First-Step Inputs}
\label{sec:data}
\label{subsec:firststep}

Table~\ref{tab:firststep_params} collects the externally assigned and first-step parameters with their values, units, and sources. The remainder of this subsection records how the Mexican data discipline the inputs that are not standard.

\begin{table}[H]
\centering
\caption{Externally assigned and first-step parameters (M\'exico calibration).}
\label{tab:firststep_params}
\begin{tabular}{llll}
\toprule
Parameter & Value & Units & Source / role \\
\midrule
\multicolumn{4}{l}{\emph{Demographics and technology}}\\
\quad Period length            & 5 & years & model frequency \\
\quad $J,\ j_R$                & $17,\ 10$ & periods & max age 100; retire at 65 \\
\quad $n_p$ (2020)            & $5.1\%$ & per 5-yr & $1.01^5-1$ \\
\quad $\alpha$                & $0.36$ & --- & capital share \\
\quad $\delta$                & $8.23\%$/yr & per 5-yr $34.9\%$ & capital depreciation \\
\quad $A$                     & $1.60$ & --- & TFP (scale) \\
\midrule
\multicolumn{4}{l}{\emph{Preferences (fixed outside SMM)}}\\
\quad $\gamma$                & $2.0$ & --- & GHH curvature \\
\quad $\nu_\ell$              & $2.0$ & --- & inverse Frisch \\
\quad $\beta$                 & $0.998^5$ & per 5-yr & discount factor \\
\midrule
\multicolumn{4}{l}{\emph{Fiscal / pension}}\\
\quad $\tau^c,\ \tau^\omega,\ \tau^k$ & $0.16,\ 0.20,\ 0.20$ & --- & statutory M\'exico rates \\
\quad $\tau^m$ (baseline), $g_y$      & $0,\ 0.19$ & --- & medical tax; $G/Y$ \\
\quad $\kappa$                & $0.50$ & --- & pension replacement rate \\
\midrule
\multicolumn{4}{l}{\emph{Labour (ENOE 2018--2024, this paper)}}\\
\quad $\theta_L,\ \theta_H$   & $-0.373,\ +0.373$ & log points & skill intercepts ($111\%$ premium) \\
\quad $e^g_j$                 & profile & efficiency units & male 20--24 $\equiv 1$; F/M $\approx0.74$ \\
\midrule
\multicolumn{4}{l}{\emph{Productivity and health (anchored)}}\\
\quad $\rho,\ \sigma_\varepsilon$ & $0.782,\ 0.265$ & --- & AR(1) of $\eta$ \\
\quad $\varrho(\theta_L),\ \varrho(\theta_H)$ & $0.30,\ 0.20$ & --- & health--productivity slope \\
\quad $\delta^h_j$            & $0.02\!\to\!0.80$ & per 5-yr & age-rising health depreciation \\
\quad $h_0$                   & $1.0$ & health units & initial health \\
\bottomrule
\end{tabular}

\vspace{0.5em}
\begin{minipage}{0.92\textwidth}
\footnotesize \textit{Notes.} Rates with ``/yr'' are annual; per-5-yr values are the model-period magnitudes used in computation. $\tau^p$ is not listed because it is endogenous (it clears the PAYG block). The six preference and health-technology parameters estimated by SMM appear separately in Table~\ref{tab:theta_hat}.
\end{minipage}
\end{table}

\paragraph{Demographics and mortality.}
The sex-specific survival schedules $\psi^g_j$ are taken from the WHO Global Health Observatory life table for M\'exico (2019), aggregated to five-year bands and closed out above age~85 by a Gompertz fit, consistent with UN WPP 2024 medium-variant magnitudes. Life expectancy at birth in M\'exico was 75.1 years in 2023 (72.1 for men, 78.3 for women), and the six-year female survival advantage is the differential the gender extension exploits. The comparative demographic and health-status landscape for the three country cases is reported in Appendix~\ref{app:supplementary}. The model lets survival depend on the endogenous health stock; a frailty--mortality probit on ENASEM 2018--2024 validates this design and disciplines the eventual mapping.\footnote{A frailty--mortality probit on ENASEM 2018--2024 (Authors' calculations) yields a pseudo-$R^2$ of $0.14$ in the full sample and $0.16$ for women, and documents the gender frailty paradox (women accumulate more health deficits yet exhibit lower mortality). Replacing the externally sourced schedule with a parametric $\psi^g_j(h)$ via this mapping is a refinement reserved for the next stage (Section~\ref{subsec:roadmap}); the present draft retains the WHO/WPP schedule.}

\paragraph{Labour productivity.}
The labour block is the team's own estimate from Mexican micro-data and is the part of the calibration that most distinguishes this paper from a literature-parameterized OLG model. Taking logs of the productivity equation~\eqref{eq:labor_productivity} gives the decomposition
\begin{equation}
\label{eq:logwage_decomp}
\log w_{i,t} \;=\; \underbrace{\log e^{g_i}_{j(i,t)}}_{\text{age-by-sex profile}} \;+\; \underbrace{\theta_{i}}_{\text{skill intercept}} \;+\; \underbrace{\varrho(\theta_i)\, h_{i,t}}_{\text{health-productivity slope}} \;+\; \eta_{i,t}.
\end{equation}
From the ENOE 2018--2024 panel we recover the sex-by-band age-efficiency profile $e^g_j$ (skill-averaged log labour income net of the skill intercept, normalized so the male 20--24 cell equals one) and the skill intercepts $\theta_g$. The male profile is concave, peaking at ages 45--49, and the female-to-male efficiency ratio averages $\approx 0.74$ at prime age---a sex earnings differential of about $26\%$. The prime-age (35--49) cross-skill log gap is $0.745$, so $\theta_L=-0.373$ and $\theta_H=+0.373$, a high-skill earnings premium of $\exp(0.745)-1\approx 111\%$. This premium is large, and it is the empirical fact that drives the sharp separation of welfare across types in the experiments. The productivity-shock process $(\rho,\sigma_\varepsilon)=(0.782,0.265)$ and the health--productivity slopes $\varrho(\theta)$ are held at upper-middle-income anchors; estimating them from the ENOE residual covariance and the self-reported-health interaction is a refinement for the next stage (Section~\ref{subsec:roadmap}). The persistence of this process matters for aggregate saving, and Section~\ref{subsec:calibrated_baseline} shows it is the dominant driver of the calibrated capital stock.

\paragraph{Health depreciation.}
The age-rising health-depreciation rate $\delta^h_j$ in~\eqref{eq:health_formation} is anchored to the force-of-aging schedule of the Grossman/deficit-accumulation tradition \citep{DalgaardStrulik2014}, rising from $2\%$ per period at young ages to near-total depreciation at the terminal age. A within-model estimator that nets medical investment out of longitudinal ENASEM/CRELES health stocks to recover $1-\delta^h_j$ directly is specified in Appendix~\ref{app:calibration} and reserved for the structural-integration stage.

\paragraph{Fiscal block.}
The tax instruments are set to statutory Mexican rates: a $16\%$ value-added tax ($\tau^c$), and labour- and capital-income tax rates of $20\%$ ($\tau^\omega,\tau^k$), with $\tau^m=0$ at baseline and a public-goods share $g_y=0.19$. The replacement rate is $\kappa=0.50$. The contribution rate $\tau^p$ is not a free parameter---it clears the PAYG block through~\eqref{eq:pension_contribution}---and the stationary debt level $B$ is the residual that closes the general budget through~\eqref{eq:gov_stationary}. Both are equilibrium objects, reported with the calibrated baseline in Section~\ref{subsec:calibrated_baseline}. The comparative fiscal and health-financing landscape for the three countries---of which the salient fact for the Mexican calibration is the high out-of-pocket share ($\approx 41\%$ of current health expenditure, against $\approx 22\%$ in Costa Rica), which leaves households more exposed to medical-expense risk---is collected in Appendix~\ref{app:supplementary}.

\subsection{Second-Step Calibration: Simulated Method of Moments}
\label{subsec:secondstep_smm}

The six parameters not pinned by external evidence or first-step regressions are estimated jointly by SMM. The estimated vector is
\[
\Theta_{\mathrm{SMM}} \;=\; \bigl(\Psi,\; \Xi,\; \xi,\; \overline{H}_0,\; h^{\mathrm{slope}},\; \zeta_h\bigr),
\]
the labour-disutility scale $\Psi$, the health-amenity level and curvature $(\Xi,\xi)$ in $s(h)$, and the three health-production parameters that enter the technology adopted following \citet{Carrasco2026},
\begin{equation}
\label{eq:healthproduction}
H_{j}(m) \;=\; \overline{H}_{j}\, m^{\zeta_{h}}, \qquad \overline{H}_{j} \;=\; \overline{H}_{0}\,\exp\bigl(-h^{\mathrm{slope}}\, j\bigr),
\end{equation}
a specific realization of the generic form in~\eqref{eq:health_formation}. The discount factor $\beta$ is pinned externally at $0.998^5$, outside the SMM loop. The six parameters are matched to six moments: prime-age male hours, the VSL, the cross-sectional income elasticity of medical spending, mean medical spending at ages 25--35, the log-slope of mean medical spending over the working life, and the within-age income elasticity of medical spending. Prime-age male hours is the Mexican (ENOE) moment; the other five are anchored to the medical-spending and VSL literature for an upper-middle-income economy.

\begin{table}[H]
\centering
\caption{Targeted moments and their sources.}
\label{tab:targeted_moments}
\begin{tabular}{llll}
\toprule
Parameter & Targeted moment & Target & Source \\
\midrule
$\Psi$               & prime-age male hours                    & $0.2686$            & ENOE 2018--2024 (this paper) \\
$\Xi$                & value of a statistical life (VSL)       & USD 2.0M            & OECD benefit transfer (Lit.) \\
$\xi$ (\emph{frozen}) & cross-section income elasticity of $m$ & $1.000$             & \citet{hall_jones_2007} (Lit.) \\
$\overline{H}_0$     & mean medical spending, ages 25--35      & $0.040$ (model units) & Hall--Jones; ENIGH (Lit.) \\
$h^{\mathrm{slope}}$ & log-slope of mean $m$, ages 25--75      & $0.120$             & Literature \\
$\zeta_h$            & within-age income elasticity of $m$     & $0.400$             & Hall--Jones; RAND (Lit.) \\
\bottomrule
\end{tabular}
\end{table}

The estimator minimizes the quadratic distance between simulated and empirical moments,
\begin{equation}
\label{eq:smm_objective}
\widehat\Theta_{\mathrm{SMM}} \;=\; \arg\min_{\Theta}\; \bigl[m^{\mathrm{sim}}(\Theta) - m^{\mathrm{dat}}\bigr]'\, W\, \bigl[m^{\mathrm{sim}}(\Theta) - m^{\mathrm{dat}}\bigr],
\end{equation}
with $W$ diagonal in the inverse squared standard errors of the data moments. Simulated moments are computed directly from the stationary distribution $\phi^g_j(\theta)$ over the four agent types weighted by birth shares, which removes Monte Carlo noise and yields a smoother objective. The six parameters are identified from the six moments, with two near-flat directions handled by holding the amenity curvature $\xi$ fixed at $0.5$ (resolving the $\Xi\leftrightarrow\xi$ ridge on the VSL target) and acknowledging an $\overline{H}_0\leftrightarrow\zeta_h$ ridge on the medical-spending moments; the full identification analysis, the Jacobian elasticity heatmap, and the multistart convergence record are in Appendix~\ref{app:calibration}.

\subsection{Calibration Results}
\label{subsec:calibration_results}

Table~\ref{tab:theta_hat} reports the six estimates with sandwich standard errors from the numerical Jacobian at $\widehat\Theta$. The labour-disutility scale $\Psi$ is sharply identified off prime-age hours ($t\approx 58$); the amenity level $\Xi$ is moderately identified. The health-production scale $\overline{H}_0$ and curvature $\zeta_h$ carry large standard errors---the acknowledged $\overline{H}_0\leftrightarrow\zeta_h$ ridge---so their point estimates are within literature ranges but imprecise; $\xi$ is frozen and its standard error is undefined.

\begin{table}[H]
\centering
\caption{Mexican-data SMM estimates $\widehat\Theta$ (best of four starts, $Q=13.28$).}
\label{tab:theta_hat}
\begin{tabular}{lrrrr}
\toprule
Parameter & Init & $\widehat\theta$ & Rel.\ change & SE \\
\midrule
$\Psi$ (labor disutility)                  & 14.00  & 13.45    & $-3.9\%$  & 0.232 \\
$\Xi$ (health amenity)                     & 0.500  & 0.275    & $-45.1\%$ & 0.171 \\
$\xi$ (amenity curvature)                  & 0.500  & 0.500    & frozen    & ---   \\
$\overline{H}_0$ (health-production scale) & 0.300  & 0.247    & $-17.5\%$ & 0.949 \\
$h^{\mathrm{slope}}$ (age-decline)         & 0.000  & $-0.0358$ & ---      & 0.138 \\
$\zeta_h$ (health-production curvature)    & 0.500  & 0.579    & $+15.8\%$ & 0.895 \\
\bottomrule
\end{tabular}

\vspace{0.5em}
\begin{minipage}{0.92\textwidth}
\footnotesize \textit{Notes.} $\widehat\theta$ is the best of four multistart Nelder--Mead runs ($Q=13.28$ at start~4; the other starts converge to $Q\approx 20.5$). $\xi$ is held fixed at $0.5$ to resolve the $\Xi\leftrightarrow\xi$ ridge. Standard errors are sandwich estimates from a numerical Jacobian at $\widehat\theta$; the wide $\overline{H}_0$ and $\zeta_h$ intervals reflect the medical-spending ridge documented in Appendix~\ref{app:calibration}.
\end{minipage}
\end{table}

Table~\ref{tab:moment_match} reports the moment fit. Five of six moments lie within two standard errors of their targets, three within one; prime-age male hours is matched essentially exactly despite its tight standard error, because $\Psi$ is nearly orthogonal to the health block. The objective consistency check $\sum_k t_k^2 = 13.28$ coincides with the recorded $Q(\widehat\theta)$. The single exception is the VSL, which diverges structurally and is the subject of Section~\ref{subsec:vsl_structural}.

\begin{table}[H]
\centering
\caption{Moment match at $\widehat\theta$.}
\label{tab:moment_match}
\begin{tabular}{lrrlrr}
\toprule
Moment                         & Data         & SE        & Source                & Model     & $t$-stat \\
\midrule
Prime-age male hours           & 0.2686       & 0.000946  & ENOE (this paper)     & 0.2687    & $+0.07$  \\
VSL (2022 USD)                 & 2{,}000{,}000 & 700{,}000 & OECD transfer (Lit.)  & 27{,}698  & $-2.82$  \\
Cross-section $\eta_m$         & 1.000        & 0.300     & Hall--Jones (Lit.)    & 0.471     & $-1.76$  \\
Mean $m$, ages 25--35          & 0.040        & 0.020     & Hall--Jones; ENIGH    & 0.0131    & $-1.35$  \\
$d\log\bar m/dj$, $j\in[2,12]$ & 0.120        & 0.050     & Literature            & 0.136     & $+0.33$  \\
Within-age $\eta_m$            & 0.400        & 0.150     & Hall--Jones; RAND     & 0.483     & $+0.56$  \\
\bottomrule
\end{tabular}

\vspace{0.5em}
\begin{minipage}{0.92\textwidth}
\footnotesize \textit{Notes.} Source column distinguishes the ENOE-this-paper moment from the retained literature/external anchors. Five of six moments are within $2$ SE (three within $1$ SE); the VSL is the structural divergence of Section~\ref{subsec:vsl_structural}. $\sum_k t_k^2 = 13.28 = Q(\widehat\theta)$.
\end{minipage}
\end{table}

\subsection{The Value of a Statistical Life}
\label{subsec:vsl_structural}

The one moment the model cannot match is the VSL: at $\widehat\theta$ the model implies USD~$27.7$k against the OECD benefit-transfer target of USD~$2.0$M. The gap is not a units artifact, and establishing that is the point of this subsection. The model VSL is the marginal rate of substitution between mortality risk and money, $\mathrm{VSL}=\beta\,\E[V_{j+1}]/u_c$, evaluated at prime age and converted to dollars by the single factor $\text{usd\_per\_unit\_c}=35{,}000$ of Section~\ref{subsec:units}; the conversion involves no period rescaling and is, if anything, conservative. The smallness is structural, and it is a property of the felicity function rather than of the amenity term. With GHH-CRRA felicity $u=(z^{1-\gamma}-1)/(1-\gamma)$ and $\gamma=2$, the flow utility of a living household is $u=1-1/z$, which the normalization sets to zero at $z=1$---exactly the value assigned to death. A bare CRRA flow whose zero coincides with the value of death mechanically delivers a near-zero willingness to pay for survival: the marginal value of being alive, relative to the marginal value of consumption, is bounded by the small per-period felicity gap between living and dying. No setting of the amenity scale $\Xi$ repairs this, because $s(h)$ shifts the \emph{level} of the composite, not the location of the felicity zero relative to death.

The fix is a separate flow value of being alive---an additive constant in felicity, as in \citet{hall_jones_2007} and the value-of-life literature---which is absent from the current specification. We read the divergence as identifying that missing object, not as a failure of fit. Two facts support treating it as a model-design item rather than a calibration target to be forced. First, the OECD transfer value is itself an income-elasticity extrapolation from US estimates; revealed-preference Mexican wage-risk studies cluster an order of magnitude lower, around USD~$200$--$325$k, so the $2.0$M target is an upper bound on the relevant Mexican VSL. Second, the VSL moment carries a standard error of USD~$700$k and is therefore nearly unweighted in the objective~\eqref{eq:smm_objective}, so the estimator correctly declines to distort the sharply identified medical-spending and hours moments to chase it. The next stage adds the flow value of being alive and re-targets the VSL against the revealed-preference Mexican range; the welfare statements in this paper, which rest on consumption and the labour wedge rather than on the monetized value of survival, do not depend on closing this gap.

\subsection{Refinements Reserved for the Next Stage}
\label{subsec:roadmap}
Three inputs now anchored to the literature will be replaced with Mexican micro-data: the survival schedule, by a parametric $\psi^g_j(h)$ from the ENASEM frailty--mortality mapping together with a GBD/DALY-based depreciation schedule and an ENASEM-estimated $\varrho(\theta)$; the medical-spending and VSL targets, by ENIGH and GBD profiles by age and household income (and a flow value of being alive, per Section~\ref{subsec:vsl_structural}); and the productivity process $(\rho,\sigma_\varepsilon)$, by the ENOE residual covariance. The estimation pipeline is unchanged under these substitutions, so the comparison between the present and the refined estimates is itself informative.

\subsection{The Calibrated Baseline Equilibrium}
\label{subsec:calibrated_baseline}

Substituting $\widehat\theta$ and the first-step inputs into the full stationary general equilibrium yields the calibrated Mexican economy of Table~\ref{tab:calibrated_baseline}---the object the experiments of Section~\ref{sec:experiments} perturb. The economy supplies effective labour $L=4.99$ and holds capital $K=4.96$ (per unit of cohort mass), producing output $Y=7.96$ at an annual interest rate of $4.21\%$. The PAYG system clears at a contribution rate $\tau^p=9.29\%$ against a system dependency ratio of $0.186$; households consume $48\%$ of output and spend $2.7\%$ on medical care.

\begin{table}[H]
\centering
\caption{Calibrated baseline equilibrium (M\'exico, 2020): stationary aggregates and welfare.}
\label{tab:calibrated_baseline}
\begin{tabular}{lc}
\toprule
Quantity & Value \\
\midrule
$K$ (capital, per cohort mass)              & 4.96 \\
$L$ (effective labor)                       & 4.99 \\
$Y$ (output)                                & 7.96 \\
$r$ (annual)                                & 4.21\% \\
$w$ (wage)                                  & 1.022 \\
$\tau^p$ (PAYG contribution rate)           & 9.29\% \\
$\overline{pen}$ (pension flow per retiree) & 0.365 \\
$C/Y$                                       & 48.1\% \\
$M/Y$                                       & 2.73\% \\
$N^R/N^W$ (PAYG dependency)                 & 0.186 \\
$5K/Y$ (annual capital--output ratio)       & 3.11 \\
$5B/Y$ (annual debt--output ratio)          & 165\% \\
\midrule
$\mathcal{W}_1(M,\theta_L)$                 & $-3.64$ \\
$\mathcal{W}_1(M,\theta_H)$                 & $-0.42$ \\
$\mathcal{W}_1(F,\theta_L)$                 & $-4.72$ \\
$\mathcal{W}_1(F,\theta_H)$                 & $-1.90$ \\
\bottomrule
\end{tabular}

\vspace{0.5em}
\begin{minipage}{0.92\textwidth}
\footnotesize \textit{Notes.} Stationary per-cohort-mass units (Section~\ref{subsec:markets}). Stock--flow ratios are annualized by the factor of five of Section~\ref{subsec:units}: $5K/Y$ and $5B/Y$ convert the five-year-flow denominator to an annual one. $r$ is annualized geometrically. Closure: $\tau^p$ clears the PAYG block; the stationary debt $B=2.62$ residually closes the general budget through~\eqref{eq:gov_stationary} (so $B/Y=0.329$ at five-year frequency, $5B/Y=165\%$ annual). Welfare $\mathcal{W}_1(g,\theta)$ is expected lifetime utility at birth ($\gamma=2$); because GHH--CRRA felicity is not homogeneous in consumption, only the \emph{ordering} across types is interpreted, not the levels (Euler residual $\max\log_{10}=-5.37$, mean $-7.42$).
\end{minipage}
\end{table}

\paragraph{Untargeted aggregates.}
None of the standard macroeconomic ratios is a calibration target: the discount factor, capital share, and depreciation rate are set externally, and the SMM moments concern hours and medical spending. The model nonetheless reproduces them. The annual capital--output ratio is $5K/Y=3.11$, in line with Mexican data; the annual return of $4.2\%$ is in the empirically plausible range; and the consumption share of $48\%$ is reasonable for an economy with a $19\%$ public-goods share and life-cycle investment. That these untargeted objects land in the right place is the validation the calibration passes. The one untargeted ratio that is off is the implied public debt: at $165\%$ of annual GDP it far exceeds M\'exico's federal debt (near $50\%$). This is a consequence of the residual debt closure---$B$ absorbs whatever stationary stock balances the primary account at the calibrated tax-and-spending configuration---rather than a targeted object, and disciplining it directly (targeting a debt ratio and letting a tax instrument clear) is a refinement for the next stage. We flag it here rather than suppress it, because the debt level matters for the infeasibility result of Section~\ref{subsec:c2_infeasibility}.

\paragraph{Lifecycle structure.}
Figure~\ref{fig:lifecycle_calibrated} plots the cohort-mean lifecycle profiles of the calibrated economy by sex and skill. Consumption and assets are hump-shaped; medical spending rises through middle age and flattens as the survival return to investment falls; and the health stock deteriorates faster for low-skill types, who carry the steeper productivity penalty and the lower earnings that shorten their effective planning horizon. These are the qualitative regularities the model was built to reproduce, and they hold at the calibrated parameters.

\begin{figure}[H]
\centering
% TODO: regenerate from run1prime_lifecycle.csv (calibrated Run 1') and upload to the project root.
\includegraphics[width=0.95\textwidth]{lifecycle_calibrated.png}
\caption{Calibrated lifecycle profiles for the four types $(g,\theta)$: consumption, labour supply, assets, medical spending, health stock, and surviving mass by age period. The hump-shape in $c$ and $a$ is the standard pattern; the divergent health trajectories by skill---high-skill households maintain health substantially longer---are the regularity that motivates the model's skill-by-health heterogeneity.}
\label{fig:lifecycle_calibrated}
\end{figure}

\paragraph{Aggregate saving is governed by income persistence.}
The calibrated economy holds less than half the capital of a symmetric-stub comparison economy---the same calibrated preferences and health technology, but with the higher AR(1) persistence and symmetric birth shares of a literature parameterization, which holds $K\approx 10.5$. Decomposing the $-53\%$ move into its two sources is itself a result about what governs aggregate saving in this class of economy. Table~\ref{tab:k_decomp} reports a four-corner experiment: each cell is a separate general-equilibrium solve at the calibrated parameters, varying only the productivity process and the birth-share vector. The Shapley decomposition attributes $-3.64$ of the $-5.53$ fall to the income process and $-1.88$ to cohort composition. Persistence dominates: lowering $\rho$ from $0.98$ to $0.782$ turns near-permanent income shocks into transitory ones, and the buffer-stock saving motive that near-unit-root risk generates collapses---even though the innovation variance rises, raising the unconditional variance from $0.063$ to $0.181$, because transitory risk is cheap to self-insure through consumption smoothing. Cohort composition works in the same direction but is smaller: M\'exico's empirical birth shares place $77.5\%$ of cohorts on the low-skill types (against $50\%$ under symmetry), and low-skill cohorts earn, save, and supply less labour. The lesson is that the aggregate capital stock of a calibrated emerging-economy life-cycle model is a statement about the persistence of idiosyncratic income risk first, and about the skill composition of the population second.

\begin{table}[H]
\centering
\caption{Aggregate capital under the income process and birth-share legs (four-corner decomposition).}
\label{tab:k_decomp}
\begin{tabular}{lcc}
\toprule
Aggregate capital $K$ & symmetric $\pi_{\rm birth}$ & asymmetric $\pi_{\rm birth}$ \\
\midrule
high persistence ($\rho=0.98$, $\sigma_\varepsilon=0.05$)   & 10.49 & 8.21 \\
data persistence ($\rho=0.782$, $\sigma_\varepsilon=0.265$) & 6.45  & \textbf{4.96} \\
\bottomrule
\end{tabular}

\vspace{0.5em}
\begin{minipage}{0.92\textwidth}
\footnotesize \textit{Notes.} Each cell is a stationary general-equilibrium solve at the calibrated SMM parameters and the baseline closure, varying only the AR(1) process and the birth-share vector. Total move $\Delta K = 4.96-10.49=-5.53$ ($-52.7\%$). The symmetric Shapley split attributes $-3.64$ to the income-process re-anchoring and $-1.88$ to the birth-share asymmetry, exhausting the move (the two channels are mildly sub-additive). The income-process leg is the larger of the two.
\end{minipage}
\end{table}

The reader should carry this economy into Section~\ref{sec:experiments}: a calibrated M\'exico in which buffer-stock saving is weak because idiosyncratic income mean-reverts, the contribution rate clears the pension system at $9.3\%$, and the four types are sharply separated by the $111\%$ skill premium. The experiments perturb exactly this object.


%% ====================================================================
%% SECTION 6 — EXPERIMENTS (at the calibrated baseline)
%% ====================================================================
\section{Experiments}
\label{sec:experiments}
\label{sec:policy}

The experiments perturb the calibrated baseline of Section~\ref{subsec:calibrated_baseline}, and they are organized around the fiscal amplification loop of Section~\ref{subsec:existence}. The two single-instrument experiments isolate the loop's two links: a pension replacement-rate cut works the direct dependency-to-wedge channel, and a health subsidy works the health-to-survival-to-dependency feedback. The cost-of-inaction experiment then stresses the loop with the projected demographic transition, and the infeasibility result shows what happens when the wedge is not allowed to absorb that stress. A forward-looking formalization design closes the section. Each stationary experiment is the calibrated baseline plus a single perturbation, solved on the same general-equilibrium code under the baseline closure ($\tau^p$ clears the PAYG block, the stationary debt $B$ residually closes the general budget); all converge cleanly and pass the model's internal diagnostics (Appendix~\ref{app:calibration}).

\subsection{Two Channels of the Amplification Loop}
\label{subsec:exp_kappa_taum}

Two experiments hold demographics at the 2020 calibration and move one policy lever each: a pension reform that lowers the replacement rate from $\kappa=0.50$ to $\kappa=0.30$, and a health subsidy that lowers the medical-spending tax from $\tau^m=0$ to $\tau^m=-0.20$. Table~\ref{tab:policy_aggregates} reports both against the calibrated baseline.

\begin{table}[H]
\centering
\caption{Single-instrument policy experiments at 2020 demographics: aggregate equilibria.}
\label{tab:policy_aggregates}
\begin{tabular}{lccc}
\toprule
 & Baseline (Run 1$'$) & $\kappa: 0.50\to0.30$ & $\tau^m: 0\to-0.20$ \\
\midrule
$K$               & 4.96    & 5.49    & 5.18   \\
$L$               & 4.99    & 5.16    & 5.04   \\
$Y$               & 7.96    & 8.44    & 8.14   \\
$r$ (annual)      & 4.21\%  & 3.79\%  & 4.00\% \\
$w$               & 1.022   & 1.047   & 1.034  \\
$\tau^p$          & 9.29\%  & \textbf{5.46\%}  & \textbf{9.62\%} \\
$\overline{pen}$  & 0.365   & 0.233   & 0.371  \\
$C/Y$             & 48.1\%  & 46.5\%  & 47.6\% \\
$M/Y$             & 2.73\%  & 2.45\%  & 3.22\% \\
$N^R/N^W$         & 0.186   & 0.182   & 0.192  \\
\bottomrule
\end{tabular}

\vspace{0.5em}
\begin{minipage}{0.92\textwidth}
\footnotesize \textit{Notes.} All three equilibria are solved on a single calibrated code path under the baseline closure; $\tau^p$ is the five-year-period contribution rate, $r$ annualized. The implied annual debt--output ratios ($5B/Y$) are $165\%$, $173\%$, and $144\%$ across the three columns---high relative to M\'exico's federal debt and a residual of the closure, as discussed in Section~\ref{subsec:calibrated_baseline}.
\end{minipage}
\end{table}

\paragraph{Pension reform ($\kappa: 0.50\to0.30$).} Cutting the replacement rate lowers the equilibrium contribution rate from $9.29\%$ to $5.46\%$, the bulk of it the direct effect of the lower $\kappa$ in the PAYG-balance identity $\tau^p=\kappa\,N^R/N^W$. The smaller public pension crowds in private retirement saving: capital rises $10.7\%$ and output $6.0\%$, compressing the annual return from $4.21\%$ to $3.79\%$ and lifting the wage. Medical spending falls slightly ($M/Y$ from $2.73\%$ to $2.45\%$) as the lower pension reduces old-age health investment, edging the dependency ratio down.

\paragraph{Health subsidy ($\tau^m: 0\to-0.20$).} The $20\%$ subsidy lowers the effective price of medical spending and raises $M/Y$ from $2.73\%$ to $3.22\%$. The induced health improvement operates through the four channels the model isolates---productivity (higher $h$ raises effective labour supply and output, $Y$ up, the labour- and capital-tax bases widen), the consumption-amenity term $s(h)$, survival, and the resource cost of $m$---and these pull the fiscal accounts in opposite directions. Through the survival channel longer lives raise the PAYG dependency ratio from $0.186$ to $0.192$, and with $\kappa$ fixed the contribution rate rises mechanically from $9.29\%$ to $9.62\%$. The subsidy is therefore not a standalone fiscal instrument: it eases pressure through the productivity channel while adding to it through the survival channel, exactly the tension the four-channel mechanism anticipates, here resolved quantitatively.

\subsection{Welfare Incidence}
\label{subsec:welfare_incidence}

We report lifetime welfare at birth $\mathcal{W}_1(g,\theta)$---expected discounted utility for a newborn of each type---as levels (Table~\ref{tab:policy_welfare}). The proper instrument for cross-type incidence is the consumption-equivalent variation (CEV) that solves, for each type,
\begin{equation}
\label{eq:cev_newborn}
\mathbb{E}_1\!\left[\sum_{j=1}^{J} \beta^{j-1}\!\left(\prod_{i=2}^{j}\psi_i^g(h_i^p)\right) u\bigl((1+\Delta_{g,\theta})\,c_j^p,\,\ell_j^p,\,h_j^p\bigr)\right] = \mathcal{W}^{p'}_1(g,\theta),
\end{equation}
the constant proportional consumption increment $\Delta_{g,\theta}$ that equates welfare under regime $p$ to that under regime $p'$. We \emph{name} the CEV as the instrument but do not compute it in this draft (see below); the incidence claims here are restricted to those robust to the welfare metric.

\begin{table}[H]
\centering
\caption{Lifetime welfare at birth $\mathcal{W}_1(g,\theta)$ (levels) under the stationary experiments.}
\label{tab:policy_welfare}
\begin{tabular}{lcccc}
\toprule
$\mathcal{W}_1(g, \theta)$ & M, $\theta_L$ & M, $\theta_H$ & F, $\theta_L$ & F, $\theta_H$ \\
\midrule
Baseline (Run 1$'$)    & $-3.64$ & $-0.42$ & $-4.72$ & $-1.90$ \\
$\kappa: 0.50\to0.30$  & $-3.78$ & $-0.36$ & $-4.99$ & $-1.96$ \\
$\tau^m: 0\to-0.20$    & $-3.62$ & $-0.39$ & $-4.72$ & $-1.88$ \\
\bottomrule
\end{tabular}

\vspace{0.5em}
\begin{minipage}{0.92\textwidth}
\footnotesize \textit{Notes.} Welfare levels from the full general-equilibrium solve. Levels are not normalized across types (the GHH felicity with an additively-separable health amenity is not homogeneous in consumption), so only sign and direction are interpreted here; magnitudes of cross-type differences require the CEV of equation~\eqref{eq:cev_newborn}.
\end{minipage}
\end{table}

Three metric-robust claims follow. First, the replacement-rate cut is \emph{regressive}: both low-skill types lose ($M,\theta_L$ and $F,\theta_L$ welfare fall), and the only type to gain is the high-skill male, who is least dependent on the public pension and benefits most from the higher equilibrium return on private saving; the high-skill female loses only mildly. The sign difference between the low-skill losses and the high-skill-male gain is robust to the welfare metric, and it is the regressive signature of the reform. Second, the health subsidy raises welfare for \emph{all four types}. We state this directional result and stop there: whether the subsidy's gains are progressive across types---larger in welfare-relevant terms for low-skill households---cannot be adjudicated from the levels, because the levels are not normalized across types; that comparison requires the CEV and is deferred. Third, demographic aging (Section~\ref{sec:aging}) reduces welfare for all four types.

The restriction to metric-robust claims is forced by the welfare metric, not chosen for caution. Because GHH--CRRA felicity is not homogeneous in consumption, the level of $\mathcal{W}_1$ is not comparable across types, and its sign---negative for the low-skill cells, near zero for the high-skill male---is an artifact of where each type's consumption composite sits relative to the felicity normalization, not a statement about welfare. Proportional comparisons of these levels would therefore be meaningless. The progressivity of the health subsidy is exactly the question the consumption-equivalent variation answers, and we report it once the welfare object is re-instrumented to compute~\eqref{eq:cev_newborn}; the directional claims above do not depend on it.

\subsection{Cost of Inaction under 2050 Demographics}
\label{sec:aging}

Holding policy at the calibrated 2020 baseline and shifting demographics to the projected 2050 structure gives the \emph{cost of inaction}: the stationary equilibrium the economy converges to when the demographic transition is fully internalized but no policy instrument adjusts. This is a long-run incidence statement, not a transition path. The demographic shift is taken from the UN World Population Prospects 2024 revision (medium variant): M\'exico's life expectancy at birth rises from $75.1$ years toward $79.8$ by 2050, fertility falls toward $1.70$, and the five-year-period population growth rate $n_p$ moves from positive to $-1.98\%$ (a contracting entering cohort). The implied demographic old-age dependency ratio roughly doubles, from $0.37$ to $0.74$, consistent with the ECLAC projection for the region \citep{ECLAC2022}; methodology and provenance for the Brass-logit survival shift are in Appendix~\ref{app:demog2050}. Within the model's cohort-weighted stationary distribution the PAYG dependency ratio $N^R/N^W$ rises from $0.186$ to $0.317$, a factor of $1.71$. We solve this under the joint closure (closure C1), in which $\tau^p$ is free to clear the PAYG block---the \emph{feasible} response. Table~\ref{tab:aging_c1} reports it against the calibrated baseline.

\begin{table}[H]
\centering
\caption{Cost of inaction under 2050 demographics (aging C1, Run 4) vs.\ the calibrated baseline.}
\label{tab:aging_c1}
\begin{tabular}{lcc}
\toprule
 & Run 1$'$ (2020) & C1 (2050, joint closure) \\
\midrule
$K$               & 4.96    & 7.13   \\
$L$               & 4.99    & 6.32   \\
$Y$               & 7.96    & 10.56  \\
$r$ (annual)      & 4.21\%  & 3.43\% \\
$w$               & 1.022   & 1.070  \\
$\tau^p$          & 9.29\%  & \textbf{15.86\%} \\
$\overline{pen}$  & 0.365   & 0.376  \\
$C/Y$             & 48.1\%  & 50.8\% \\
$M/Y$             & 2.73\%  & 2.81\% \\
$N^R/N^W$         & 0.186   & 0.317  \\
\bottomrule
\end{tabular}

\vspace{0.5em}
\begin{minipage}{0.92\textwidth}
\footnotesize \textit{Notes.} Baseline closure C1: $\tau^p$ clears the PAYG block at the 2050 demographics. $\tau^p$ at five-year-period frequency, $r$ annualized; annual $5B/Y=132\%$ under C1.
\end{minipage}
\end{table}

The headline is the contribution rate: under unchanged policy, $\tau^p$ rises from $9.29\%$ to $15.86\%$---a $71\%$ increase that tracks the rise in the dependency ratio through the PAYG-balance condition $\tau^p=\kappa\,N^R/N^W$. The economy accumulates capital ($K$ up $43.8\%$) as the strengthened life-cycle saving motive at longer life expectancy outweighs the contracting cohort, lifting output and compressing the annual return from $4.21\%$ to $3.43\%$. The lesson is that inaction is not free: even with the contribution rate left to clear the pension system, the demographic shock alone pushes the labour-tax wedge up by roughly two-thirds, with the welfare cost (Section~\ref{subsec:welfare_incidence}) falling on every type.

\subsection{The Feasible Policy Menu: Non-Existence under Debt-Financed Inaction}
\label{subsec:c2_infeasibility}

The aging response above is feasible because the contribution rate is allowed to clear the pension system. The conceptual hinge of the paper is what happens when it is not. Consider instead the closure that holds both the contribution rate $\tau^p$ and the benefit formula at their 2020 values and absorbs the resulting PAYG deficit with general government debt---``do nothing on the pension side and borrow.'' Under this closure the calibrated economy has no finite stationary equilibrium at M\'exico's 2050 demographics.

\begin{proposition}[Non-existence under debt-financed inaction]
\label{prop:c2}
At the calibrated 2050 demographics, with $\tau^p$ and the benefit formula pinned at their 2020 values and the PAYG deficit absorbed by debt, no finite stationary general equilibrium exists. The pension deficit forces the stationary debt stock so high that the net return is driven below the population growth rate, and the stationary-debt relation that pins debt at a finite level breaks down.
\end{proposition}

The mechanism is transparent in the stationary-debt relation $B=\text{primary}/(r^n-n_p)$, which fixes the debt stock that services itself out of the primary balance only while the net return $r^n$ exceeds the population growth rate $n_p$. The pinned contribution rate leaves a growing pension deficit, pushing the primary balance into deficit and the implied debt stock upward; financing it requires ever more capital, so the marginal product of capital and hence $r^n$ fall. In the iteration record (Figure~\ref{fig:c2_crash}), capital climbs to $K\approx 44.9$ and $r^n$ crosses $n_p$ from above at the eleventh step, with the pension deficit at $2.57$ per cohort mass; as the denominator $r^n-n_p$ passes through zero the implied debt ricochets from $-49$ to $+170$ in a single step and the solver no longer converges. The crossing, not any level, is the result.

This is economics, not a numerical artifact. The shortfall pushes the economy into the dynamic-inefficiency region $r^n<n_p$, where a permanent primary deficit cannot be rolled over against a finite stationary debt: the transversality condition that bounds debt fails. Something must give---the contribution rate (closure C1 above), the benefit formula, or general taxation. The menu of feasible long-run responses to demographic aging is therefore narrower than the policy debate assumes: pure debt-financing of the pension shortfall is structurally unavailable, not merely costly. The result holds at the calibrated baseline and is robust to pinning the contribution rate at either the calibrated or the higher pre-harmonization 2020 rate, since any fixed rate below the clearing rate leaves a deficit that grows with the demographic transition.\footnote{The iteration in Figure~\ref{fig:c2_crash} pins $\tau^p$ at $10.83\%$; pinning it at the calibrated $9.29\%$ enlarges the deficit and reaches the $r^n<n_p$ crossing sooner. The non-existence is generic to the closure, not an artifact of the pinned value.}

\begin{figure}[H]
\centering
% TODO: regenerate c2_interp1_crash_panel.png from the calibrated Run 0 trace
% (run0_history.csv: crossing at iteration 11, K -> 44.9, deficit 2.57, B -49 -> +170)
% and upload to the Overleaf project root. The current panel reflects the pre-calibration trace.
\includegraphics[width=\textwidth]{c2_interp1_crash_panel.png}
\caption{Failure to admit a finite stationary equilibrium under the debt-financed-inaction closure. \textit{Left:} aggregate capital across outer general-equilibrium iterations, rising monotonically and jumping at the crossing. \textit{Middle:} the net interest rate $r^n$ and the population growth rate $n_p$ ($-0.40\%$ annual); $r^n$ crosses $n_p$ from above at iteration~11, entering the dynamic-inefficiency region (shaded), where the stationary-debt relation becomes ill-conditioned. \textit{Right:} the goods-market residual across iterations, oscillating without convergence---in contrast to the geometric decay of the baseline solve.}
\label{fig:c2_crash}
\end{figure}

\subsection{Formalization Experiment}
\label{sec:formalization}

The experiments above study how fiscal instruments absorb demographic pressure within the formal-sector economy. A complementary experiment, which we design here and quantify in the next stage, asks the fiscal consequences of expanding the contributory base itself.

Let the working-age population be partitioned into formal workers $\mathcal{F}$ and informal workers $\mathcal{I}$, with high-skill workers overrepresented in $\mathcal{F}$ and low-skill workers in $\mathcal{I}$---the equilibrium outcome of contribution costs, benefit valuations, enforcement, and credibility documented in the formalization literature \citep{WorldBankInformalizing2025, BoschMelguizoPages2013}. A formalization program transfers a fraction $\alpha\in(0,1]$ of $\mathcal{I}$ into the contributory system; the newly formalized workers enter with their observed ENOE income, health, and survival profiles, and the model is re-solved for the new stationary equilibrium. Two opposing fiscal forces result. The \emph{contribution-base effect} raises the mass of contributors $N^W$ (retirees rise only with a lag, as the newly formalized must accumulate contribution years), lowering the dependency ratio and the equilibrium $\tau^p$ and widening the labour-tax base. The \emph{expenditure-cost effect} works the other way: the newly formalized are predominantly low-skill, with earlier health depreciation, lower productivity, and lower lifetime earnings, raising per-capita health claims and---through the survival channel---potential pension receipt. The model quantifies, for each country and formalization rate $\alpha$, the change in $\tau^p$, the welfare incidence on incumbent formal workers, and the threshold $\alpha^*$ at which the contribution-base effect offsets the expenditure-cost effect; a targeted health subsidy or a reduced pension vesting threshold can be layered on to study complementarities, connecting the analysis directly to the reform roadmap of \citet{WorldBankInformalizing2025}.

\subsection{Caveats}
\label{subsec:exp_caveats}

Four caveats qualify the experiments. First, the analysis compares stationary equilibria; the within-transition incidence on cohorts alive in 2020, and the general-taxation closure in which broad taxation rather than the contribution rate absorbs the pension shortfall, are left for the transition analysis. Second, cross-type welfare incidence---the progressivity of the health subsidy---awaits the consumption-equivalent variation of~\eqref{eq:cev_newborn}. Third, the implied stationary debt level is high and is a residual of the closure rather than a disciplined object (Section~\ref{subsec:calibrated_baseline}); targeting a debt ratio is a natural next refinement. Fourth, the experiments inherit the inputs still held at literature anchors---the productivity process, the depreciation schedule, and the health--productivity penalty---whose Mexican-data refinement is described in Section~\ref{subsec:roadmap}.


%% ====================================================================
%% SECTION 8 — CONCLUSION
%% ====================================================================
\section{Conclusion}
\label{sec:conclusion}

This paper builds a stochastic overlapping-generations model with continuous endogenous health capital, heterogeneous agents who differ by sex and skill, endogenous factor prices, and a pay-as-you-go pension system whose contribution rate clears endogenously, and calibrates it to M\'exico. The sex-by-age labour-efficiency profile, the skill premium, and the prime-age hours target are estimated from the ENOE 2018--2024 panel; the remaining six preference and health-technology parameters are estimated by simulated method of moments. The calibration matches five of six target moments within two standard errors and three within one. The one moment it does not match, the value of a statistical life, is informative rather than a failure of fit: a GHH--CRRA felicity whose zero coincides with the value of death cannot generate a large willingness to pay for survival, which identifies a missing flow value of being alive as the next modeling step (Section~\ref{subsec:vsl_structural}). Three substantive results stand out. The strongest is structural: at M\'exico's projected 2050 demographics, debt-financed inaction on the pension side admits no finite stationary equilibrium (Proposition~\ref{prop:c2}), so the menu of feasible long-run responses is narrower than the policy debate assumes. The second is that the aggregate capital stock of the calibrated economy is governed first by the persistence of idiosyncratic income risk and only second by skill composition---the empirical income process mean-reverts, and buffer-stock saving is correspondingly weak. The third is that a replacement-rate cut is regressive, while the equilibrium contribution rate rises by roughly seventy percent under unchanged policy at 2050 demographics.

The model's central analytical contribution is the identification of a fiscal amplification loop that links health investment to pension sustainability through the longevity channel. Health investment raises the health stock, which improves survival probabilities and extends working lives---strengthening the life-cycle savings motive, raising aggregate capital, and compressing the equilibrium interest rate. But longer lives also raise the dependency ratio $N^R/N^W$, mechanically increasing the pension contribution rate $\tau^p = \kappa \cdot N^R/N^W$ and widening the tax wedge on formal labor income. The net fiscal effect of improved population health is therefore ambiguous in sign. This ambiguity cannot be resolved by partial-equilibrium analysis: it requires the calibrated general-equilibrium model that this paper builds, in which health-driven changes in aggregate savings move factor prices and feed back into every household's decision problem.

The paper develops two complementary experimental frameworks. The policy counterfactual battery (Section~\ref{subsec:exp_kappa_taum}) studies how fiscal instruments---pension parameters ($\kappa$, $j_R$), tax rates ($\tau^\omega$, $\tau^c$, $\tau^k$), and health subsidies ($\tau^m$)---can absorb the fiscal pressures generated by demographic change. The four-channel mechanism through which health subsidies operate---productivity, utility, survival, and resource cost---makes $\tau^m$ qualitatively different from the other fiscal instruments and motivates its joint analysis with pension reforms. The formalization experiment (Section~\ref{sec:formalization}) studies the fiscal consequences of expanding the contributory base by incorporating informal workers with their observed income and health profiles, identifying the threshold formalization rate $\alpha^*$ at which the contribution-base effect dominates the expenditure-cost effect. Both frameworks are designed to evaluate policy complementarities---whether combinations of instruments dominate any single-instrument response---which is the most policy-relevant and computationally demanding question the model addresses.

The model's architecture supports a natural extension toward gender-differentiated fiscal analysis. By expanding the agent-type space from two skill types to four---$(\theta_L^M, \theta_H^M, \theta_L^F, \theta_H^F)$---calibrated with sex-specific earnings trajectories and the survival differentials documented in the data section, the model can decompose fiscal sustainability pressures by sex within the existing equilibrium framework. This calibration-based approach lets observed income profiles and survival probabilities carry the information that a structural household bargaining model would generate endogenously, at a fraction of the computational cost.

The refinements described in Section~\ref{subsec:roadmap}---the ENASEM frailty--mortality mapping and a parametric $\psi^g_j(h)$, ENIGH/GBD medical-spending targets and a flow value of being alive, the productivity process from ENOE residuals, and the consumption-equivalent variation for cross-type welfare incidence---will sharpen the estimates and extend the analysis to Costa Rica and Panam\'a, alongside the within-transition incidence and the formalization experiment.

The principle that organizes the analysis is straightforward: the question facing Latin American health systems is not whether to spend more, but when to spend, through which institutional channels, and with what consequences for the contributory architecture on which fiscal sustainability depends.


%% ====================================================================
%% APPENDICES
%% ====================================================================
\begin{appendices}

\section{Fiscal Gap Methodology}
\label{app:fiscal_gap}

The fiscal gap for country $c$ in year $t$ is defined as:
\begin{equation}
\label{eq:fiscal_gap}
G_{c,t} = H_{c,t}^{\text{est}} - H_{c,t}^{\text{actual}}
\end{equation}
where $H_{c,t}^{\text{est}}$ denotes estimated health expenditure requirements and $H_{c,t}^{\text{actual}}$ denotes actual public health expenditure, both expressed as shares of GDP. Estimated requirements are constructed by applying a base-year age-expenditure profile $h(j)$ to the projected population structure:
\begin{equation}
\label{eq:fiscal_gap_est}
H_{c,t}^{\text{est}} = \frac{\sum_j h(j) \cdot N_{j,c,t}}{Y_{c,t}}
\end{equation}
where $N_{j,c,t}$ is the population of age group $j$ in country $c$ at time $t$, and $Y_{c,t}$ is nominal GDP. The base-year expenditure profile $h(j)$ is estimated from national health accounts microdata for the reference year.


\section{2050 Demographic Inputs: Derivation and Sources}
\label{app:demog2050}

\paragraph{Source.} The demographic primitives for the comparative-steady-state experiment of Section~\ref{sec:aging} are anchored to the UN World Population Prospects (WPP) 2024 revision, online edition (accessed 24 May 2026), M\'exico medium-variant projections. The aggregate anchors are: life expectancy at birth in 2050 of 79.8 years (sex gap of 5.6 years, implying male life expectancy of 77.0 and female of 82.6); a total fertility rate of 1.70; and a total population peaking near 2042 at 152 million, declining to 150.6 million by 2050. The 2020 baseline reference is PAHO (2024), reporting M\'exico life expectancy in 2023 of 75.1 years (male 72.1, female 78.3).

\paragraph{Method.} The aggregate WPP targets do not directly fix the 17-band, sex-specific survival schedule required by the model. We therefore apply a one-parameter Brass-logit shift to the 2020 baseline pooled survival schedule,
\begin{equation}
\label{eq:brass_logit}
\operatorname{logit}\bigl(1 - p_{2050, g, j}\bigr) \;=\; \alpha_g + \operatorname{logit}\bigl(1 - p_{2020, j}\bigr),
\end{equation}
with $\alpha_g$ chosen separately for males ($\alpha_m = -0.158$) and females ($\alpha_f = -0.696$) to hit the sex-specific WPP life-expectancy targets. The population growth rate $n_p$ is set to $-0.4\%$ annual (a five-year-period rate of $-1.98\%$), consistent with the implied entering-cohort contraction under the WPP 2024 total fertility rate. Together with the Brass-shifted survival schedules, this yields a stationary old-age dependency ratio of 0.74, exactly twice the baseline value of 0.37, matching the ECLAC anchor for M\'exico by 2045 \citep{ECLAC2022}.

\paragraph{Caveat.} The Brass shift preserves the shape of age-conditional mortality and adjusts its level through a single parameter per sex. It cannot capture compositional changes in mortality, such as disproportionate gains at older ages from chronic-disease management. For the comparative-steady-state purpose pursued here this is adequate; for the September three-country deliverable the placeholder is replaced with WPP single-age life tables aggregated to five-year bands, and the gender decomposition is re-run on calibrated rather than stub primitives.

\section{Calibration and Computational Appendix}
\label{app:calibration}

This appendix documents the apparatus behind the Mexican-data calibration of Section~\ref{sec:calibration}: the SMM objective and weighting, the identification analysis, the multistart convergence record, the parity and Euler diagnostics, the override mechanics that let the calibration run without editing the protected solver, and the first-step provenance.

\subsection{SMM Objective and Weighting}
\label{app:smm_objective}

The estimator minimizes the quadratic form~\eqref{eq:smm_objective}, $Q(\Theta)=[m^{\rm sim}(\Theta)-m^{\rm dat}]'\,W\,[m^{\rm sim}(\Theta)-m^{\rm dat}]$, with $W$ diagonal in the inverse squared standard errors of the six data moments, so that each moment contributes its squared $t$-statistic to the objective and $Q(\widehat\Theta)=\sum_k t_k^2$. The weights are therefore the moment precisions reported in Table~\ref{tab:moment_match}; the tightly estimated prime-age hours moment (SE $0.000946$) is heavily weighted, while the VSL moment (SE USD $700$k) is nearly unweighted, which is why the optimiser is willing to sacrifice the VSL fit (Section~\ref{subsec:vsl_structural}). At each candidate $\Theta$ the model is solved as a partial equilibrium at prices anchored to a stationary general-equilibrium solve, and simulated moments are computed directly from the stationary distribution $\phi^g_j(\theta)$ over the four agent types weighted by birth shares---removing the Monte Carlo noise a panel simulation would introduce and yielding a smoother objective for the simplex. Asymptotic standard errors on $\widehat\Theta$ are sandwich estimates from a numerical Jacobian of $m^{\rm sim}$ at $\widehat\Theta$; because $\Xi$ enters $s(h)$ inside the GHH composite, its identification interacts with the externally fixed $\gamma$, which is treated as conditioning information.

\subsection{Identification}
\label{app:identification}

The pre-optimization Jacobian at the starting values (Figure~\ref{fig:jacobian_elast}, reported as target-relative elasticities) shows that every free parameter has strong sensitivity on at least one moment, with no globally flat direction. Two near-flat directions require handling. The first is the $\Xi\leftrightarrow\xi$ ridge on the VSL target: both the amenity level and its curvature move \texttt{vsl\_usd} with the same sign and comparable magnitude, so they cannot be separately identified from that single moment. We resolve it by freezing $\xi=0.5$ during estimation, which makes the system over-identified in the $\xi$ dimension. The second is the $\overline{H}_0\leftrightarrow\zeta_h$ ridge on the medical-spending moments: the scale and curvature of the health-production function compete with opposite signs, $\zeta_h$ roughly three times the stronger, so their point estimates are within literature ranges but imprecise (the wide standard errors in Table~\ref{tab:theta_hat}). This ridge is acknowledged rather than resolved; it does not prevent identification of $\overline{H}_0$ to within range. One benign off-diagonal sign-flip (the within-age elasticity with respect to $h^{\rm slope}$) appears at the boundary $h^{\rm slope}=0$, where the sign is fragile to the income profile; it is not a new identification problem.

\begin{figure}[H]
\centering
% TODO: ensure 02b_jacobian_elasticity.png is uploaded to the Overleaf project root
% (source: Calibration/outputs/diagnostics/02b_jacobian_elasticity.png; 02_jacobian.png is the alternative).
\includegraphics[width=0.85\textwidth]{02b_jacobian_elasticity.png}
\caption{Pre-optimization Jacobian: target-relative elasticities of the six moments with respect to the six parameters. The $\Xi\leftrightarrow\xi$ ridge on \texttt{vsl\_usd} (resolved by freezing $\xi$) and the strong $\zeta_h$ column on the medical-spending moments (the acknowledged $\overline{H}_0\leftrightarrow\zeta_h$ ridge) are the two identification features behind the estimation restrictions of Section~\ref{subsec:secondstep_smm}.}
\label{fig:jacobian_elast}
\end{figure}

\subsection{Multistart Protocol and Convergence}
\label{app:multistart}

The objective is minimized by a multistart Nelder--Mead simplex: four starts of at most forty iterations each, run under four-way contention at roughly twelve minutes per partial-equilibrium evaluation. The smoke test at the starting values returns $Q(\theta_{\rm init})=73.11$ with all six moments finite. The best of the four starts is start~4 at $Q=13.28$; the remaining starts converge to $Q\approx 20.5$, so the reported optimum is well separated from the runner-up basins. Figure~\ref{fig:obj_trace} traces the running-best objective across evaluations on a log scale.

\begin{figure}[H]
\centering
% TODO: regenerate 03_objective_trace.png for the Mexican-data multistart (best Q = 13.28)
% and upload to the Overleaf project root
% (source: Calibration/outputs/diagnostics/03_objective_trace.png reflects the intermediate run).
\includegraphics[width=0.9\textwidth]{03_objective_trace.png}
\caption{Multistart progress: running-best objective $Q$ versus evaluation index, on a log scale. The descent terminates at the best-of-four optimum $Q=13.28$.}
\label{fig:obj_trace}
\end{figure}

\subsection{Parity and Euler Diagnostics}
\label{app:parity_euler}

The SMM scaffold reproduces the production \texttt{Household-Gender} solver to machine precision: the maximum absolute cohort-mean difference between the two solvers is below $10^{-10}$ (achieved at $7\times10^{-11}$), so the calibration and the experiments share an identical household solution algorithm. At the calibrated baseline (Run~1$'$) the Euler-equation residual has a maximum of $\log_{10}=-5.37$ and a mean of $-7.42$, with the mass-weighted bulk of the distribution comfortably below the convergence gate; the two-sex welfare-gap diagnostic also passes. The calibrated experiments (Runs~2--4) each converged within twenty-six outer iterations and pass the same internal gates.

\subsection{Override Mechanics}
\label{app:overrides}

The calibration and the experiments ran without editing the protected solver \texttt{ge\_model\_gender.jl}. Three override patterns were used, each respecting Julia's compilation semantics. The AR(1) persistence and innovation scale $(\rho,\sigma_\varepsilon)$ are \texttt{const} scalars inlined into the Rouwenhorst discretization, so they were passed as \emph{literals} into the model initializer rather than redefined (a \texttt{const} redefinition would silently fail to propagate into already-compiled methods). The birth-share vector $\pi_{\rm birth}$ is a \texttt{const} array, so in-place mutation propagates, as do the ENOE age-efficiency and skill-intercept overrides. The 2050 demographics are installed on typed globals (population growth and the sex-specific survival schedules) through a setter rather than by inlining. A frozen-parameter bug in the diagnostics module was worked around through a shim rather than a source edit. These patterns are recorded for replicability.

\subsection{First-Step Provenance}
\label{app:provenance}

The survival schedule is the WHO Global Health Observatory life table for M\'exico (2019), aggregated to five-year bands and closed out above age~85 by a Gompertz fit, consistent with UN WPP 2024 medium-variant magnitudes. The ENOE age-efficiency profile $e^g_j$ is constructed, for each sex and five-year band, as the skill-averaged log labour income net of the skill intercept, normalised so the male 20--24 cell equals one; the skill intercepts follow from the prime-age (35--49) cross-skill log gap of $0.745$. The hours target is the equal-weighted mean of the ENOE \texttt{hrsocup} share over male prime-age (30--44) skill cells. Unit conversions for the retained VSL and medical-spending moments, the income-box recipe for $e^g_j$, and the full source attribution for every input are documented in \texttt{inputs\_mxdata/provenance.md}.

\section{Supplementary Tables and Data}
\label{app:supplementary}

This appendix collects the three-country comparative landscape that motivates the project. Only M\'exico is calibrated in the present draft; these indicators discipline none of the parameters in Section~\ref{sec:calibration} and are reported for context.

\subsection{Demographic and Health-Status Indicators}
\label{app:demog_table}

\begin{table}[H]
\centering
\caption{Demographic and Health Status Indicators}
\label{tab:demog}
\begin{tabular}{lccc}
\hline\hline
\textbf{Indicator} & \textbf{Mexico} & \textbf{Costa Rica} & \textbf{Panama} \\
\hline
Life expectancy at birth, total (2023)           & 75.1 & 81.0 & 78.1 \\
Life expectancy, male (2023)                     & 72.1 & 78.6 & 75.4 \\
Life expectancy, female (2023)                   & 78.3 & 83.5 & 80.9 \\
Fertility rate (births per woman, 2023)          & 1.82 & 1.51 & 2.30 \\
Infant mortality (per 1{,}000 live births, 2022) & 12.4 & 8.7  & 14.5 \\
Population aged 65+ (\% of total, 2023)          & 8.1  & 10.4 & 8.9  \\
Diabetes prevalence (\% of adults)               & 16.9 & 9.8  & 11.0 \\
\hline\hline
\end{tabular}
\begin{minipage}{\linewidth}
\smallskip
\footnotesize
\textit{Sources:} \citet{PAHOHIA2024}; \citet{WHOGHO2024};
\citet{WorldBankWDI2024}; \citet{OECDHaG2023}.
\end{minipage}
\end{table}

\subsection{Health Expenditure Indicators}
\label{app:health_table}

\begin{table}[H]
\centering
\caption{Health Expenditure Indicators (Reference Year: 2022)}
\label{tab:health_exp}
\begin{tabular}{lccc}
\hline\hline
\textbf{Indicator} & \textbf{Mexico} & \textbf{Costa Rica} & \textbf{Panama} \\
\hline
Current health expenditure (\% of GDP)       & 5.5          & 7.1   & 7.8   \\
Public health expenditure (\% of CHE)        & $\approx$51  & $\approx$69 & $\approx$72 \\
Out-of-pocket expenditure (\% of CHE)        & $\approx$41$^{a}$ & 22.4 & $\approx$26 \\
OOP per capita, USD current                  & $\approx$253 & 220   & $\approx$290 \\
Total health exp.\ per capita, USD PPP       & $\approx$1{,}100 & 1{,}658 & $\approx$1{,}450 \\
OOP as \% of household final consumption     & $\approx$5.0 & 2.6   & $\approx$3.5 \\
\hline\hline
\end{tabular}
\begin{minipage}{\linewidth}
\smallskip
\footnotesize
\textit{Sources:} \citet{WHOGHED2024}; \citet{OECDHaG2023};
\citet{WHOGHO2024}. CHE = Current Health Expenditure; OOP =
Out-of-pocket; PPP = Purchasing Power Parity. Values marked
($\approx$) are approximations. $^{a}$~Mexico's OOP share
corresponds to the 2021 data vintage, the most recent available.
\end{minipage}
\end{table}

\subsection{Fiscal Indicators}
\label{app:fiscal_table}

\begin{table}[H]
\centering
\caption{Fiscal Indicators}
\label{tab:fiscal}
\begin{tabular}{lccc}
\hline\hline
\textbf{Indicator} & \textbf{Mexico} & \textbf{Costa Rica} & \textbf{Panama} \\
\hline
Standard VAT rate                                  & 16\% (IVA) & 13\% (IVA) & 7\% (ITBMS) \\
Top personal income tax rate (2023)                & 35\%  & 25\%  & 25\%  \\
Social security contributions (employer + employee) & $\approx$30\% & $\approx$37\% & $\approx$19\% \\
General government revenue (\% of GDP, 2022)       & $\approx$23 & $\approx$24 & $\approx$21 \\
Tax revenue (\% of GDP, 2022)                      & $\approx$17 & $\approx$14 & $\approx$12 \\
\hline\hline
\end{tabular}
\begin{minipage}{\linewidth}
\smallskip
\footnotesize
\textit{Sources:} \citet{PwCWWTS2024}; \citet{OECDRevStats2024};
\citet{WorldBankWDI2024}; \citet{CCSS2023}.
\end{minipage}
\end{table}

\subsection{Labour Market Indicators}
\label{app:labour_table}

\begin{table}[H]
\centering
\caption{Labour Market Indicators (Reference Year: 2022)}
\label{tab:labour}
\begin{tabular}{lccc}
\hline\hline
\textbf{Indicator} & \textbf{Mexico} & \textbf{Costa Rica} & \textbf{Panama} \\
\hline
Male labour force participation rate (\%)    & 76.7        & 72.4        & 77.6        \\
Female labour force participation rate (\%)  & 45.4        & 50.7        & 51.8        \\
Gender participation gap (male/female ratio) & 1.69        & 1.43        & 1.50        \\
Gender wage gap (\% of male wage)            & $\approx$14 & $\approx$18 & $\approx$19 \\
Informality rate (\% of employed)            & $\approx$55 & $\approx$43 & $\approx$47 \\
Unemployment rate (\%)                       & 3.3         & 11.2        & 7.7         \\
Gini coefficient                             & 0.43        & 0.48        & 0.49        \\
\hline\hline
\end{tabular}
\begin{minipage}{\linewidth}
\smallskip
\footnotesize
\textit{Sources:} \citet{ILOILOSTAT2023}; \citet{ECLACILO2023};
\citet{IMF2023}; \citet{WEFGGR2024}; \citet{WorldBankWDI2024};
\citet{WorldBankPovCalNet2022}.
\end{minipage}
\end{table}

\subsection{Life Table for Mexico}
\label{app:survival}

The following life table corresponds to Mexico and is used to
calibrate the age-specific survival probabilities $\psi_{j}$ in the
model. Data are drawn from \citet{UISP2019}.

\begin{table}[H]
\centering
\caption{Life Table for Mexico: Mortality and Survival Probabilities}
\label{tab:survival_mx}
\begin{tabular}{lcc}
\hline\hline
\textbf{Age group} & \textbf{Probability of death} & \textbf{Survival probability} \\
\hline
$<$ 1 year   & 0.011576 & 0.9884 \\
1--4         & 0.002282 & 0.9977 \\
5--9         & 0.001132 & 0.9989 \\
10--14       & 0.001470 & 0.9985 \\
15--19       & 0.003858 & 0.9961 \\
20--24       & 0.005897 & 0.9941 \\
25--29       & 0.006550 & 0.9935 \\
30--34       & 0.007966 & 0.9920 \\
35--39       & 0.011217 & 0.9888 \\
40--44       & 0.013710 & 0.9863 \\
45--49       & 0.019216 & 0.9808 \\
50--54       & 0.028413 & 0.9716 \\
55--59       & 0.042568 & 0.9574 \\
60--64       & 0.063162 & 0.9368 \\
65--69       & 0.093160 & 0.9068 \\
70--74       & 0.137294 & 0.8627 \\
75--79       & 0.208926 & 0.7911 \\
80--84       & 0.312093 & 0.6879 \\
85+          & 1.000000 & 0.0000 \\
\hline\hline
\end{tabular}
\begin{minipage}{\linewidth}
\smallskip
\footnotesize
\textit{Source:} \citet{UISP2019}. Net probability of death and
survival probability by age group for the total population of Mexico.
\end{minipage}
\end{table}

\end{appendices}


%% ====================================================================
%% BIBLIOGRAPHY
%% ====================================================================
\bibliography{references}

\end{document}
